\PassOptionsToPackage{unicode}{hyperref}
\PassOptionsToPackage{hyphens}{url}
\PassOptionsToPackage{dvipsnames,svgnames,x11names}{xcolor}
\documentclass[
  11pt,
]{article}
\usepackage{amsmath,amssymb}
\usepackage{iftex}
\ifPDFTeX
  \usepackage[T1]{fontenc}
  \usepackage[utf8]{inputenc}
  \usepackage{textcomp} % provide euro and other symbols
\else % if luatex or xetex
  \usepackage{unicode-math} % this also loads fontspec
  \defaultfontfeatures{Scale=MatchLowercase}
  \defaultfontfeatures[\rmfamily]{Ligatures=TeX,Scale=1}
\fi
\usepackage{lmodern}
\ifPDFTeX\else
  % xetex/luatex font selection
  \setmainfont[]{Latin Modern Roman}
  \setsansfont[]{DejaVu Sans}
  \setmonofont[]{DejaVu Sans Mono}
  \setmathfont[]{Latin Modern Math}
\fi
\IfFileExists{upquote.sty}{\usepackage{upquote}}{}
\IfFileExists{microtype.sty}{% use microtype if available
  \usepackage[]{microtype}
  \UseMicrotypeSet[protrusion]{basicmath} % disable protrusion for tt fonts
}{}
\makeatletter
\@ifundefined{KOMAClassName}{% if non-KOMA class
  \IfFileExists{parskip.sty}{%
    \usepackage{parskip}
  }{% else
    \setlength{\parindent}{0pt}
    \setlength{\parskip}{6pt plus 2pt minus 1pt}}
}{% if KOMA class
  \KOMAoptions{parskip=half}}
\makeatother
\usepackage{xcolor}
\usepackage[letterpaper,margin=0.84in]{geometry}
\usepackage{longtable,booktabs,array}
\usepackage{calc} % for calculating minipage widths
\usepackage{etoolbox}
\makeatletter
\patchcmd\longtable{\par}{\if@noskipsec\mbox{}\fi\par}{}{}
\makeatother
\IfFileExists{footnotehyper.sty}{\usepackage{footnotehyper}}{\usepackage{footnote}}
\makesavenoteenv{longtable}
\setlength{\emergencystretch}{3em} % prevent overfull lines
\providecommand{\tightlist}{%
  \setlength{\itemsep}{0pt}\setlength{\parskip}{0pt}}
\setcounter{secnumdepth}{5}
\ifLuaTeX
\usepackage[bidi=basic]{babel}
\else
\usepackage[bidi=default]{babel}
\fi
\babelprovide[main,import]{american}
\ifPDFTeX
\else
\babelfont{rm}[]{Latin Modern Roman}
\fi
\let\LanguageShortHands\languageshorthands
\def\languageshorthands#1{}

\usepackage{fancyhdr}

\usepackage{microtype}
\usepackage{longtable,booktabs,array}
\definecolor{Ink}{HTML}{173345}
\definecolor{Accent}{HTML}{176B77}
\pagestyle{fancy}
\fancyhf{}
\setlength{\headheight}{15pt}
\setlength{\emergencystretch}{3em}
\setlength{\parskip}{5pt}
\setlength{\parindent}{0pt}
\allowdisplaybreaks[2]
\ifLuaTeX
  \usepackage{selnolig}  % disable illegal ligatures
\fi
\IfFileExists{bookmark.sty}{\usepackage{bookmark}}{\usepackage{hyperref}}
\IfFileExists{xurl.sty}{\usepackage{xurl}}{} % add URL line breaks if available
\urlstyle{same}



\usepackage{etoolbox}
\makeatletter
\providecommand{\subtitle}[1]{% add subtitle to \maketitle
  \apptocmd{\@title}{\par {\large #1 \par}}{}{}
}
\makeatother

\title{Projected cross terms and compact retained-source jets}\author{}\date{}
\providecommand{\diver}{\operatorname{div}}
\begin{document}\maketitle
\hypertarget{part-i.-actual-retained-source-energy-and-projected-cross-terms}{%
\section{Part I. Actual retained-source energy and projected cross
terms}\label{part-i.-actual-retained-source-energy-and-projected-cross-terms}}

This note does not prove the missing extensive
low-spectrum ordinary bound. It gives a new self-contained marked
reduction for a concrete realization B\_* of the alternative
original-law observations described in the compact Wilson boundary calculation.
It keeps every live normalizer in the actual law, avoids differentiation
of sharp component fibers, and controls the projected ordinary/current
cross term by an exact dual-Dirichlet identity. B\_* is not identified
with an unspecified retained B, the observation X of~\cite{CT}, or a native R-II variable.
Below P\_B denotes this explicit P\_\{B\_*\}.

The scalar separated-marginal reduction, equation~(5) of~\cite[\S1]{OS}, is not
reproved here. The result here is an
independent extension to the random rare-component/safe-seed
observation, under the geometric properties stated explicitly below. No
new proof of that scalar formula is claimed.

\hypertarget{fixed-original-law-source-and-local-hypotheses}{%
\subsection{Fixed original law, source, and local
hypotheses}\label{fixed-original-law-source-and-local-hypotheses}}

Let mu be the complete original finite-volume Wilson law on the product
of compact link groups. Use beta \textgreater= 1 and elementary
plaquettes with distinct links. Its density with respect to product Haar
measure is strictly positive. At original links e let K\_e be
conditional expectation given all other original links, D\_e=I-K\_e, and
define

\[
 A=\sum_e D_e,\qquad
 \mathcal E_A(v)=\langle v,Av\rangle
                 =\sum_e\|D_ev\|_2^2.
 \tag{1.1}
\]

Each K\_e is the actual normalized Wilson heat-bath kernel. The identity

\[
 \mathcal E_A(v)=\frac12\sum_e
       \mathbb E|v(U^e)-v(U)|^2
 \tag{1.2}
\]

uses the stationary coupling in which U has law mu and U\^{}e is
obtained by resampling only link e from that kernel. Both marginals are
exactly mu. No approximation, good-field truncation, or tensorization
across different e is involved.

Retain the complete original ordinary source

\[
 F_f=R_f=\sum_c f_c r_c^0,
 \qquad M_f=\sum_c |f_c|^2.
 \tag{1.3}
\]

The roots c include the original fixed pointing/averaging incidence. All
original root coordinates and color multiplicities are counted once. The
marks f\_c are deterministic. The actual physical multiplier
c\_beta\^{}\{-1\} a\^{}\{-2\} is restored only at the end.

The original same-Wilson fixed-record moment premise implies, for the
complete literal local row,

\[
 \sup_c\|r_c^0\|_{L^q(\mu)}\le C q
 \quad\text{for every fixed finite }q.
 \tag{1.4}
\]

Indeed the force-square part is bounded by C beta times the sum of
plaquette actions in a fixed determining neighborhood, and the full
local Casimir row is bounded. This uses the complete row, including
neighboring and pointed contributions; it does not assume a moment bound
for a deleted component of a row.

Only q=4 is needed for the new Dirichlet estimate. The linear q growth
in (1.4) is for the original ordinary row, from its local exponential
moment. An additional local source grade may instead have a stated bound
C\_\{q,grade\} (for example polynomial growth q\^{}k); that dependence
must remain. The Dirichlet proof still uses its actual L4 bound, while
the rare-residual logarithm changes with its moment growth. No uniform
bound across an unbounded collection of inherited grades is asserted.

The complete total-current row satisfies

\[
 T_e=\Delta_eS_0-\beta|\nabla_eS_0|^2,
 \qquad K_eT_e=0,
 \qquad \|T_e\|_2^2\le C_T.
 \tag{1.5}
\]

The last bound holds unconditionally by~\cite{CT}. The source identity is

\[
 h_f=R_f+\tau_f,\qquad
 \tau_f=\frac13\sum_e\widetilde f_e T_e,
 \qquad\sum_e|\widetilde f_e|^2\le C_{\rm inc}M_f.
 \tag{1.6}
\]

The tilde only records the actual fixed pointing/averaging incidence if
the source roots c and link marks e are listed differently. It is the
exact original incidence map, not a newly chosen source. Reference~\cite{CT} gives

\[
 \mathbb E\tau_f=0,
 \qquad \|\tau_f\|_2^2\le C M_f.
 \tag{1.7}
\]

\hypertarget{the-exact-observation-and-geometric-locality-needed}{%
\subsection{The exact observation and geometric locality
needed}\label{the-exact-observation-and-geometric-locality-needed}}

Let D be the exact set of bad plaquettes, defined by s\_p \textgreater{}
epsilon\^{}2. Let the finite-range collection of factor supports include
all original action factors, all good and bad selector tests, and every
complete source determining set/pointed offset. A source determining set
is included in this graph whether or not it is a density factor. This
changes only the component grouping; it does not change the Wilson
density.

Make a fixed graph on original plaquettes by declaring p adjacent to q
whenever some factor support intersects BOTH link boundaries of p and q.
This graph has bounded degree. Let C range over connected components of
the induced bad set D and define

\[
 E_D=\bigcup_{p\in D}\partial p,
 \qquad E_C=\bigcup_{p\in C}\partial p,
 \qquad O=(D,U\text{ outside }E_D).
 \tag{2.0}
\]

If a factor touches E\_C and E\_\{C'\}, it intersects the boundaries of
some p in C and q in C', which are adjacent by definition. Hence C=C'.
In particular the E\_C are disjoint, and a complete source readout
touches at most one rare component. Shared eliminated links also join
their plaquettes by this graph. A graph joining only bad plaquettes
through geometric face adjacency would not provide this factorization
statement.

On the original links, join two links when a factor reads both, and
choose a deterministic independent seed set S\_0 of positive density in
this bounded-degree graph. Thus every action/test/readout factor reads
at most one seed. Retain only the safe seeds S(D), namely those s in
S\_0 for which every factor containing s avoids E\_D. Define the
concrete observation

\[
 \mathcal B_*=(D,U\text{ outside }E_D\cup S(D)),\qquad
 P_B=\mathbb E(\,\cdot\mid\mathcal B_*).
 \tag{2.1}
\]

No action/test/readout factor touches both a safe seed and E\_D. Every
plaquette containing a safe seed is good: a bad such plaquette would
place the seed itself in E\_D and violate safety. Conditional on
B\_\emph{, the complete selected law therefore factors over the rare
components and safe seeds. Its factors are the original compact selected
component fibers and the exact sharp good-star fibers, with their actual
normalizers Z\_C and Z\_\{s,good\}. These normalizers remain in the
marginal density of B\_}.

For a normalized expectation of one local root r\_c\^{}0, factors from
fibers that the root does not read cancel exactly between numerator and
denominator. This cancellation is permissible only inside that genuine
normalized conditional expectation. It does not remove the corresponding
Z factors from the actual marginal density, nor make them constant under
a later resampling.

The construction has the following locality properties, independently of
any one-link spectral gap. Changing one original link e changes bad
flags only for plaquettes containing e, and changes E\_D only in their
fixed neighborhood. Connected component labels can nevertheless merge or
split along a whole old or new component. Each C has a fixed-thickness
determining boundary N(C), formed by all original factors touching E\_C.
The conditional mean m\_c=P\_B r\_c\^{}0 is the literal normalized
integral over the component or safe seed read by r\_c\^{}0, if any. It
depends only on that fiber, its boundary, the local selector labels and
the source's remaining variables.

The original multipoint tail can be derived from the exact positive
transfer operator of the odd-torus transfer theorem (\texttt{local\_contacts/odd\_torus\_transfer\_moment.tex}) and its partition-power input, as follows. This
derivation strengthens its one-plaquette application; it does not infer
a multipoint estimate from the one-plaquette marginal alone.

\hypertarget{lemma-2.1-arbitrary-set-original-moments-including-odd-tori}{%
\subsubsection{Lemma 2.1: arbitrary-set original moments, including odd
tori}\label{lemma-2.1-arbitrary-set-original-moments-including-odd-tori}}

Let C consist of plaquettes of one fixed orientation, and let 0
\textless= t \textless{} 1. On the same original n\^{}4 periodic Wilson
torus, n \textgreater= 3,

\[
 \mathbb E\exp\!\left(t\beta\sum_{p\in C}s_p\right)
 \le\left(\frac{Z_{n,(1-t)\beta}}{Z_{n,\beta}}\right)^{|C|/n^4}.
 \tag{2.2}
\]

Thus, in the original size/coupling regime of the odd-torus transfer theorem, its partition-power bound
supplies a constant K\_t independent of n and beta such that this
expectation is at most K\_t\^{}\{\textbar C\textbar\}.

Proof. The generalized Schatten trace Hölder inequality is

\[
 \left|\operatorname{Tr}\prod_{i=0}^{n-1}A_i\right|
 \le\prod_{i=0}^{n-1}
        \bigl(\operatorname{Tr}A_i^n\bigr)^{1/n}
 \tag{2.3}
\]

for positive trace-class A\_i. It has no parity restriction. Take the
marked orientation to be (1,2). For time direction 3, split C into its
arbitrary slice patterns C\_i and let D\_i be multiplication by exp(t
beta sum\_\{p in C\_i\}s\_p) on that spatial slice. The positive
transfer operator T gives the EXACT original numerator

\[
 \operatorname{Tr}(D_0T D_1T\cdots D_{n-1}T)
 =\operatorname{Tr}\prod_{i=0}^{n-1}
               (T^{1/2}D_iT^{1/2}).
 \tag{2.4}
\]

Each factor in parentheses is positive and trace class. Its n-th trace
power is the exact partition function with the particular pattern C\_i
repeated in every direction-3 slice. Dividing (2.3) by the ORIGINAL Z=Tr
T\^{}n expresses the bound as a product of n same-denominator ratios,
each to power 1/n.~Empty patterns have ratio one.

Apply the same argument in direction 4 to each repeated-pattern
numerator. One now has branches indexed by original direction-3 and
direction-4 coordinates, with exponent 1/n\^{}2, whose marked sets
repeat in both transverse directions.

Use time direction 1 next. Each branch has an arbitrary list of marked
temporal slabs. In slab i the marked temporal convolution coefficients
are (1-t) beta and the unmarked ones are beta. The exact slab operator
is the positive transfer operator T\_\{lambda\_i\}=M
Q\_\{lambda\_i\}P\_gauge M. The numerator is Tr product\_i
T\_\{lambda\_i\}, so (2.3) repeats each particular slab in every time-1
position, again dividing only by the same original Z. Finally repeat in
time direction 2. Each final branch is either the complete marked
(1,2)-orientation on the torus, or empty.

Counting the original marked plaquettes through these four repetitions
gives exactly exponent \textbar C\textbar/n\^{}4 for the partition ratio
with all (1,2) couplings reduced to (1-t) beta. Reducing the other
nonnegative plaquette couplings to the same value only increases that
numerator, proving (2.2). Reversing the in-plane orientation when
changing the time direction merely relabels slab/link indices; the
original periodic integrations and mark count remain exact. All Haar
factors, action constants, and gauge projections are the ones in the odd-torus transfer theorem.
No new probability measure is used as a denominator at an intermediate
step.

For an arbitrary plaquette set C with six orientation classes, ordinary
Hölder with all six exponents equal to six gives

\[
 \boxed{\quad
 \mathbb E\exp\!\left(\frac{t\beta}{6}
                      \sum_{p\in C}s_p\right)
 \le K_t^{|C|/6}.
 \quad}
 \tag{2.5}
\]

Therefore

\[
 \boxed{\quad
 \mathbb P(C\subset D)\le e^{-\lambda |C|},\qquad
 \lambda=\frac{t\beta\epsilon^2-\log K_t}{6}.
 \quad}
 \tag{2.6}
\]

This proves the exact exponent stated in the companion notes, with the
number six arising from the orientation Hölder step. The remaining
input is the partition-power estimate of the odd-torus transfer theorem in its stated
original regime, not an assumed arbitrary-set probability tail.

These properties imply the needed damage event. If m\_c(U) differs from
m\_c(U\^{}e) and c is outside a fixed neighborhood of e, then an OLD or
a NEW rare component connects N\_0(e) to the determining neighborhood of
c.~A merger is covered by a component of D(U\^{}e); a split is covered
by a component of D(U). A safe seed can change its status only next to
an affected component. Its extra readout radius is fixed and is absorbed
in the neighborhoods.

This assertion compares the canonical fiber-integral versions of m\_c at
two configurations. At boundary values where a conditional normalizer is
zero, choose any version; such values are null for the stationary
marginal used here. The proof never asks for a uniform lower bound on
Z\_C or Z\_\{s,good\}.

\hypertarget{an-original-law-stabilization-lemma}{%
\subsection{An original-law stabilization
lemma}\label{an-original-law-stabilization-lemma}}

\hypertarget{lemma-3.1}{%
\subsubsection{Lemma 3.1}\label{lemma-3.1}}

Let m\_c be functions on the original Wilson probability space with
\textbar\textbar m\_c\textbar\textbar\_4 \textless= L, and suppose for
the stationary one-link resampling coupling

\[
 \mathbb P\{m_c(U^e)\ne m_c(U)\}
 \le A_0 e^{-\kappa d(e,c)}
 \tag{3.1}
\]

outside a fixed near neighborhood. If the kernel exp(-kappa d(e,c)/4)
has uniformly bounded row and column sums, then

\[
 \boxed{\quad
 \mathcal E_A\left(\sum_c f_c m_c\right)
 \le C L^2 M_f.
 \quad}
 \tag{3.2}
\]

The constant depends on A\_0, kappa, graph incidence and the fixed near
radius, not the lattice volume.

Proof. Put delta\_e m\_c=m\_c(U\^{}e)-m\_c(U). Stationarity gives
\textbar\textbar delta\_e m\_c\textbar\textbar\_4 \textless= 2L. The
difference is supported on the event in (3.1), so Hölder gives

\[
 \|\delta_e m_c\|_2
 \le 2L\,\mathbb P\{\delta_e m_c\ne0\}^{1/4}
 \le C L k(e,c),
 \tag{3.3}
\]

where k is constant on the fixed near neighborhood and has the stated
exponential tail elsewhere. Minkowski gives

\[
 \left\|\delta_e\sum_c f_cm_c\right\|_2
 \le C L\sum_c |f_c|k(e,c).
 \tag{3.4}
\]

Apply the row/column Schur bound to this deterministic positive kernel:

\[
 \sum_e\left(\sum_c|f_c|k(e,c)\right)^2
 \le\left(\sup_e\sum_ck(e,c)\right)
      \left(\sup_c\sum_ek(e,c)\right)M_f.
 \tag{3.5}
\]

Equation (1.2) proves (3.2). No independence between different damage
events, source roots, or heat-bath steps is assumed.

\hypertarget{application-to-the-exact-o-and-b-observations}{%
\subsubsection{Application to the exact O and B
observations}\label{application-to-the-exact-o-and-b-observations}}

For m\_c=P r\_c\^{}0, conditional Jensen gives
\textbar\textbar m\_c\textbar\textbar\_4 \textless= C\_4. If the fixed
graph has at most a\_0\^{}m connected animals of size m through a
specified vertex, the geometric implication in Section 2 and
stationarity give

\[
 \mathbb P\{m_c(U^e)\ne m_c(U)\}
 \le C\sum_{m\ge c_0d(e,c)-C_0}a_0^m e^{-\lambda m}
 \le C'e^{-\kappa d(e,c)}.
 \tag{3.6}
\]

The OLD and NEW contributions each use the same original marginal mu; no
comparison of two conditional component laws is needed. For sufficiently
large lambda the constants can be chosen uniformly. On the fixed-range
four-dimensional lattice the exponential kernel is summable. A
bounded-degree-only version also works once its exponent exceeds the
graph's growth rate.

Consequently,

\[
 \boxed{\quad
 \mathcal E_A(P_{\mathcal B}R_f)\le C M_f.
 \quad}
 \tag{3.7}
\]

The same proof applies to O with no safe seeds. This is an estimate for
the retained mean itself in the ORIGINAL heat-bath Dirichlet form. It is
stronger than a bound on the internal rare residual. It is still an
energy estimate, not an L2 susceptibility estimate.

The argument is insensitive to a cusp of a sharp-star normalizer,
because it never differentiates a sharp fiber or its moving boundary.
The normalizer's finite change is inside the complete conditional mean
m\_c, whose L4 norm follows from Jensen. Integrating another color would
change the observation; this proof does not reset its density or iterate
a Wilson product law.

\hypertarget{compressed-exact-dynamics-and-the-currents-dual-dirichlet-bound}{%
\subsection{Compressed exact dynamics and the current's dual-Dirichlet
bound}\label{compressed-exact-dynamics-and-the-currents-dual-dirichlet-bound}}

First let P be conditional expectation onto ANY deterministic
observation sigma-algebra of the original U. Let H\_P=ran P and define
the compression

\[
 A_P=(PAP)|_{H_P}.
 \tag{4.1}
\]

This is a bounded nonnegative self-adjoint operator at finite volume.
Its quadratic form on H\_P is exactly the original form:

\[
 \langle v,A_Pv\rangle=\mathcal E_A(v),\qquad v\in H_P.
 \tag{4.2}
\]

It is an actual reversible observed dynamics. An individual update
samples a complete configuration from its actual conditional law given
the observation, makes one original Wilson link update, and reobserves.
Thus this dynamics updates geometry/history labels when the original
update changes them. It is generally different from a native RG dynamics
that holds those labels fixed.

Put t\_f=P tau\_f.~For any v in H\_P, the exact score cancellation
K\_eT\_e=0 gives

\[
 \langle v,t_f\rangle
 =\frac13\sum_e\widetilde f_e\langle D_ev,T_e\rangle.
 \tag{4.3}
\]

Hence

\[
 \boxed{\quad
 |\langle v,t_f\rangle|
 \le C\sqrt{M_f\,\mathcal E_A(v)}.
 \quad}
 \tag{4.4}
\]

This holds for any such P, including deterministic observations with
sharp selectors and varying-dimensional fibers. No derivative of P or of
the observed density is used. In spectral language,

\[
 \|A_P^{-1/2}t_f\|_2^2\le CM_f,
 \qquad
 \|1_{[0,\epsilon]}(A_P)t_f\|_2^2\le C\epsilon M_f.
 \tag{4.5}
\]

The generalized inverse is on the complement of the kernel; (4.4)
implies t\_f is orthogonal to that kernel. Formula (4.5) follows from
the quadratic-form dual norm or directly from the spectral theorem. It
does not require a uniform spectral gap.

For the finite-volume original Wilson law, ker A consists only of
constants: D\_ev=0 for all e makes v invariant under every coordinate
resampling, and the strictly positive density has the same null sets as
product Haar measure. Therefore, for a deterministic original-U
observation, ker A\_P on H\_P also consists only of constants. This
assertion is specific to that compressed original dynamics. A native
generator that freezes a history label may have many additional zero
modes.

For a stochastic history observation, first realize the completed kernel
as Y=Y(U,V) on the exact product extension mu(dU) nu(dV), where V is
independent auxiliary randomness, and let A resample U at fixed V. The
proof of (4.3)--(4.5) and the projected cross bound remains valid.
However ker A now consists of functions of V; its compressed kernel can
contain nonconstant modes. For the particular centered ORIGINAL source
r\_f=P(R\_f(U)-ER\_f) with deterministic marks, these modes are
orthogonal to r\_f: if k(V) belongs to ran P, then
\textless r\_f,k\textgreater=\textless R\_f-ER\_f,k\textgreater=0 by
independence. This is a separate orthogonality argument, not a claim
that the kernel is constants. For general saved/history-dependent
coefficients, reference rows or shared-history observables, the source
reduction must retain
\textbar\textbar1\_\{\{0\}\}(A\_P)r\_f\textbar\textbar\^{}2 unless the
analogous orthogonality is proved. Neither a native frozen-history
kernel nor the stochastic lift may be silently replaced by the
constant-only kernel of deterministic B\_*.

\hypertarget{exact-projected-cross-terms}{%
\subsection{Exact projected cross
terms}\label{exact-projected-cross-terms}}

Let r\_f=P(R\_f-ER\_f), t\_f=P tau\_f, and let Q be any spectral
projection of A\_P. Since Q is self-adjoint and idempotent,

\[
 \langle Qr_f,Qt_g\rangle
 =\frac13\sum_e\widetilde g_e
      \langle D_eQr_f,T_e\rangle.
 \tag{5.1}
\]

Spectral projection contracts the compressed form, so

\[
 \boxed{\quad
 |\langle Qr_f,Qt_g\rangle|
 \le C\sqrt{M_g\,\mathcal E_A(r_f)}.
 \quad}
 \tag{5.2}
\]

For P=P\_B satisfying (3.7), this is the extensive mixed-mark bound

\[
 |\langle Qr_f,Qt_g\rangle|
 \le C\sqrt{M_fM_g}.
 \tag{5.3}
\]

In particular, the actual projected squared-norm comparison follows from
the exact expansion

\[
 \|QP(h_f-Eh_f)\|_2^2-\|Qr_f\|_2^2
 =2\langle Qr_f,Qt_f\rangle+\|Qt_f\|_2^2.
 \tag{5.4}
\]

The cross term is bounded by (5.2), and the last term by~\cite{CT}. Thus

\[
 \boxed{\quad
 \left|\|QP(h_f-Eh_f)\|_2^2-\|Qr_f\|_2^2\right|
 \le CM_f
 \quad(P=P_B).
 \quad}
 \tag{5.5}
\]

This conclusion is not obtained by squaring an error norm. The separate
exact cross-term calculation is essential. Its mixed bilinear version,
with f and g, is bounded by C sqrt(M\_fM\_g).

The same result holds for contractive spectral multipliers, with the
evident replacement of Q by q(A\_P), because the mixed pairing then has
q(A\_P)\^{}2 acting on r and its form norm still contracts.

\hypertarget{qualification-for-actual-native-output-projections}{%
\subsubsection{Qualification for actual native output
projections}\label{qualification-for-actual-native-output-projections}}

If Q is an arbitrary orthogonal projection on H\_P, (5.1) still holds,
but the valid bound is

\[
 |\langle Qr_f,Qt_g\rangle|
 \le C\sqrt{M_g\,\mathcal E_A(Qr_f)}.
 \tag{5.6}
\]

Its right side need not be bounded by the energy of r\_f.~To use a
native spectral projection Q\_native, one must prove

\[
 \mathcal E_A(Q_{\rm native}v)\le C_Q\mathcal E_A(v)
 \tag{5.7}
\]

on the actual relevant source family, or a suitable two-sided form
comparison with the native generator. This includes all native scores,
reference terms and moving-wall/surface contributions needed for that
comparison. No such comparison is proved here.

A finite-dimensional example explains why this condition cannot be
omitted. Let orthonormal u,v satisfy A u=delta u and A v=v, take
r=delta\^{}\{-1/2\}u and t=v. Then E\_A(r)=1,
\textbar\textbar t\textbar\textbar=1, and \textless r,t\textgreater=0.
Projection Q onto (u+v)/sqrt(2) gives \textless Qr,Qt\textgreater=1/(2
sqrt(delta)), which diverges as delta tends to zero. This is a
counterexample to the abstract inference that arbitrary orthogonal
projection preserves the cross bound, not a counterexample to the
desired Wilson estimate.

\hypertarget{complete-residuals-and-a-marked-low-spectrum-reduction-for-b_}{%
\subsection{Complete residuals and a marked low-spectrum reduction for
B\_*}\label{complete-residuals-and-a-marked-low-spectrum-reduction-for-b_}}

The exact construction permits independent proofs of the two complete
residual estimates. They use the SAME original source and normalization
as the Dirichlet estimate.

\hypertarget{rare-component-residual}{%
\subsubsection{Rare-component residual}\label{rare-component-residual}}

For an exact component C, let I(C) be the set of complete source roots
whose determining sets intersect E\_C, and put F\_C=sum\_\{c in I(C)\}
f\_c r\_c\^{}0. Every root belongs to at most one such component, and
\textbar I(C)\textbar{} \textless= C \textbar C\textbar. Conditional on
O, different F\_C depend on independent component fibers, so

\[
 \mathbb E\operatorname{Var}(F_f\mid O)
 =\mathbb E\sum_{C\text{ actual}}
                   \operatorname{Var}(F_C\mid O).
 \tag{6.0a}
\]

For a fixed possible C of size m, weighted Cauchy gives

\[
 |F_C|^2\le
       \left(\sum_{c\in I(C)}|f_c|^2\right)
       \left(\sum_{c\in I(C)}|r_c^0|^2\right).
 \tag{6.0b}
\]

The local exponential source moment gives the genuine original
rare-weighted moment

\[
 \mathbb E[1_E|r_c^0|^2]
 \le C p\log^2(e/p),\qquad p=\mathbb P(E).
 \tag{6.0c}
\]

One can obtain (6.0c) by Hölder with exponent of order log(e/p), using
\textbar\textbar r\_c\^{}0\textbar\textbar\_q \textless= Cq. In
particular, on the exact-component event with p \textless= exp(-lambda
m), and lambda large,

\[
 \mathbb E[1_{\{C\text{ actual}\}}|r_c^0|^2]
 \le C\lambda^2m^2 e^{-\lambda m}.
 \tag{6.0d}
\]

Equations (6.0a)--(6.0d), followed by summation over animals touching a
fixed source determining set, imply

\[
 \|(I-P_O)R_f\|_2^2
 \le C\lambda^2\sum_c|f_c|^2
        \sum_{m\ge1}m^3a_0^m e^{-\lambda m}
 \le C\lambda^2e^{-\lambda}M_f.
 \tag{6.0e}
\]

The last constant is uniform once lambda is larger than a fixed graph
threshold. Every expectation here is under the original normalized law.
In (6.0a) the conditional components have their actual Z\_C, and the
outer geometry law retains the same Z\_C. There is no independent
selected-mass factor to insert again.

\hypertarget{safe-seed-residual}{%
\subsubsection{Safe-seed residual}\label{safe-seed-residual}}

For a safe seed s, let I(s) be the complete source roots reading s and
let F\_s=sum\_\{c in I(s)\} f\_c r\_c\^{}0. The independent-seed
construction gives \textbar I(s)\textbar{} \textless= C and at most one
seed per source determining set. Safety implies that every component
contribution and its normalizer avoid all these seeds. Consequently P\_O
F\_f, conditional on B\_\emph{, is a seed-independent term plus sum\_s
F\_s. The safe-star kernels are independent conditional on B\_}, so

\[
 \|(P_O-P_B)R_f\|_2^2
 =\mathbb E\sum_{s\in S(D)}
       \operatorname{Var}(F_s\mid\mathcal B_*)
 \le C M_f.
 \tag{6.0f}
\]

For the last inequality, apply finite-cardinality Cauchy to F\_s and
then the unconditional original L2 bound for each complete r\_c\^{}0.
The indicator that s is safe is at most one. No starwise uniform
conditional moment bound, no spectral gap, and no boundary
differentiation are needed.

If a seed is unsafe, a bad plaquette lies within a fixed factor
neighborhood of it. The one-plaquette original exponential moment thus
gives

\[
 \mathbb E\#(S_0\setminus S(D))
 \le C|S_0|e^{-\gamma\beta}
 \tag{6.0g}
\]

for fixed epsilon and sufficiently large beta. This verifies that the
good-information layer has positive density. This count is not used as a
substitute for the variance calculation (6.0f).

We have therefore proved the complete residual estimates for the exact
nested observations B\_* subset O:

\[
 \|(I-P_O)R_f\|_2^2\le C\lambda^2e^{-\lambda}M_f,
 \qquad
 \|(P_O-P_B)R_f\|_2^2\le CM_f.
 \tag{6.1}
\]

Because P\_B P\_O=P\_B, these two errors are orthogonal, and

\[
 \|(I-P_B)R_f\|_2^2\le CM_f.
 \tag{6.2}
\]

Keep h=R+tau exactly. The complete residual of h has norm squared at
most CM\_f by (6.2), (1.7), and conditional contraction. Orthogonality
of conditional expectation gives

\[
 \operatorname{Var}(h_f)
 =\|(I-P_B)h_f\|_2^2+
   \|P_B(h_f-Eh_f)\|_2^2.
 \tag{6.3}
\]

Use the exact cross estimate (5.4) with Q=I to replace the second term
by \textbar\textbar r\_f\textbar\textbar\^{}2 with error CM\_f.~Finally,

\[
 \|1_{(1,\infty)}(A_B)r_f\|_2^2
 \le\langle r_f,A_Br_f\rangle\le CM_f.
 \tag{6.4}
\]

There is no nonconstant compressed zero mode. Therefore

\[
 \boxed{\quad
 \left|\operatorname{Var}_\mu(h_f)
 -\left\|1_{(0,1]}(A_B)
     \{P_BR_f-ER_f\}\right\|_2^2\right|
 \le CM_f.
 \quad}
 \tag{6.5}
\]

More generally the high-spectrum cutoff 1 can be replaced by
eta\textgreater0, at the cost CM\_f(1+eta\^{}\{-1\}) in this reduction.
The missing bound remains the low-spectrum L2 mass, not the weighted
spectral first moment.

The two-mark version is

\[
 \left|\operatorname{Cov}_\mu(h_f,h_g)
 -\left\langle1_{(0,1]}(A_B)r_f,
                 1_{(0,1]}(A_B)r_g\right\rangle\right|
 \le C\sqrt{M_fM_g}.
 \tag{6.6}
\]

For (6.6), use the orthogonal conditional decomposition for the complete
residuals, then (5.3) for both cross terms and the spectral energy bound
for both high-frequency terms. No cancellation of two uncontrolled
squared norms is assumed.

For an unnormalized complete parent measure of mass M, with the original
Wilson marginal after normalization, deterministic physical/saved
coefficients multiply (6.5)--(6.6) by their actual squared or mixed
products and by M once. In particular the physical squared multiplier is
c\_beta\textsuperscript{\{-2\}a}\{-4\}. A coefficient fixed only after
selecting a history must remain inside that history sum or have a
verified envelope. No separate future selector mask may be attached to
the already formed error without a new estimate.

\hypertarget{what-this-achieves-and-what-it-leaves}{%
\subsection{What this achieves and what it
leaves}\label{what-this-achieves-and-what-it-leaves}}

The new estimate (3.7) controls how the retained ordinary mean changes
under one ORIGINAL-law resampling. It pays for geometry changes and
component mergers/splits by original cluster probabilities and pays for
local source values by original moments. It handles sharp good-star
cusps without pretending their derivatives are bounded, and it preserves
the live induced normalizers.

The exact dual-Dirichlet identity (4.4) then controls the
ordinary/current projected cross term for the compressed actual
dynamics. Together with the already formed residual bounds, this
supplies a self-contained marked version of the low-spectrum
reduction of~\cite[\S2]{OS} for the specified distinct B observation.

It does not prove

\[
 \left\|1_{(0,1]}(A_B)(P_BR_f-ER_f)\right\|_2^2\le CM_f.
 \tag{7.1}
\]

The proved estimate is only

\[
 \int\lambda\,d\sigma_{r_f}(\lambda)\le CM_f;
 \tag{7.2}
\]

the missing quantity is the unweighted mass of that spectral measure
near zero. A field can have bounded Dirichlet energy and a large
low-spectrum L2 norm. Thus (3.7) cannot be used as a Poincare inequality
with an implicit constant one.

The complete current has the additional negative-order estimate (4.5);
the ordinary retained mean does not presently have one. Obtaining such a
source-specific estimate, or a sufficiently integrable annealed
semigroup covariance decay for r\_f, is a concrete remaining spatial
mechanism.

There is a useful exact connection back to the estimate required in~\cite{CT}. If
(7.1) were proved, then (6.4) and (6.2) would give Var\_mu(R\_f)
\textless= CM\_f.~Consequently, for the actual P\_X of~\cite{CT} and ANY subsequent
completed positive observation Y and orthogonal output Q,

\[
 \|Q\,\mathbb E[(P_XR_f-ER_f)\mid Y]\|_2^2
 \le\operatorname{Var}_\mu(R_f)\le CM_f.
 \tag{7.3}
\]

This implication uses successive conditional contractions and needs no
identification of B\_* with X or R-II and no screening relation. A
native form comparison is needed to transfer the projected CROSS-TERM
estimate before such a global variance bound is available; it is not an
additional hypothesis for the implication (7.3). Actual signed localized
representatives and saved source grades must still be identified with
their complete original pushforwards before applying it.

Similarly, if both X and B\_* have proved complete residual budgets for
the SAME original R\_f, their unprojected retained covariance forms
differ by an extensive amount through separate orthogonal
decompositions:

\[
 \operatorname{Cov}(P_XR_f,P_XR_g)
 -\operatorname{Cov}(P_BR_f,P_BR_g)
 =\langle(I-P_B)R_f,(I-P_B)R_g\rangle
  -\langle(I-P_X)R_f,(I-P_X)R_g\rangle.
 \tag{7.4}
\]

This is bounded by C sqrt(M\_fM\_g). It is not a statement about
separately projected future observations, where the orthogonality in
(7.4) can be lost.

Nothing here identifies B with the actual R-II observation, proves
comparison with a native generator that freezes histories, supplies
mixed-domain simultaneous-wall admission, or closes S94. Those remain
distinct obligations.

\hypertarget{other-routes-examined-and-not-used-as-proof}{%
\subsection{Other routes examined and not used as
proof}\label{other-routes-examined-and-not-used-as-proof}}

A directional-action integration-by-parts identity involves the TOTAL
Wilson force dotted with a directional force. It does not equal the
square of that directional force; the cross-direction terms remain. An
action specific-heat bound alone consequently does not control this
ordinary pressure direction.

At the free Gaussian level, the ordinary quadratic symbol can be
represented by a bounded nonlocal multiplier applied to the total force,
leading to an extensive bound. A nonlinear replacement requires the
actual field-dependent operator, its divergence, derivatives, and
curvature remainder. Pointwise small plaquette curvature does not by
itself give uniform positivity relative to a massless Laplacian at
arbitrarily small momenta. No such replacement is assumed in this note.

In particular, a proposed vector field V=(D/\textbar F\textbar\^{}2)F is
singular at backgrounds with total Yang--Mills force zero but nonzero
directional forces. A pointwise domination D \textless=
C\textbar F\textbar\^{}2 plus curvature terms, even when available in a
restricted chart, is not by itself a bound on a centered variance.

\hypertarget{an-exact-original-covariance-identity}{%
\subsection{An exact original covariance
identity}\label{an-exact-original-covariance-identity}}

There is a useful finite-range identity before any observation is taken.
If the determining set of the local ordinary root \(r_c^0\) does not
contain original link \(e\), then

\[
 \mathbb E(r_c^0T_e)=\mathbb E\{r_c^0 K_eT_e\}=0.
\]

Together with the same identity for two noninteracting current roots,
this implies that the covariance kernels of the complete original rows
\(h_c\) and \(R_c\) agree exactly beyond a fixed determining-stencil
distance. Their difference is a finite-range kernel, with uniformly
bounded entries under the original local moment premise. Consequently

\[
 |\operatorname{Cov}(h_f,h_g)-\operatorname{Cov}(R_f,R_g)|
 \le C\sqrt{M_fM_g}.
\]

This is a direct cross-term proof for the original rows. It does not
require inferring a variance difference from a norm error. Arbitrary
observations can alter the spatial pattern of the retained covariance,
as Part III shows; the specified compressed-generator estimate in Part I
is what controls its projected cross term.

For completeness, only the one-link second moment from~\cite{CT} is
needed in Part I. Compact integration by parts gives \(\mathbb ET_e=0\)
and

\[
 \mathbb ET_e^2=-\mathbb E\langle\nabla_eS_0,\nabla_eT_e\rangle.
\]

The Casimir relation and compact one-link Hessian bound imply
\(|\langle\nabla_eS_0,\nabla_eT_e\rangle|\le C\beta|\nabla_eS_0|^2\).
The latter expectation equals \(\mathbb E\Delta_eS_0\le6C_FN\), giving
\(\|T_e\|_2^2\le C\). Exact covariance locality and bounded incidence
then give \(\|\tau_f\|_2^2\le CM_f\). No global spectral gap is used.

\newpage

\hypertarget{part-ii.-exact-compact-marginal-source-jets}{%
\section{Part II. Exact compact marginal source
jets}\label{part-ii.-exact-compact-marginal-source-jets}}

\hypertarget{status-and-purpose}{%
\subsection{Status and purpose}\label{status-and-purpose}}

This note proves a local statement about the \textbf{actual original-law
retained ordinary source} at the deterministic separated-link
observation used in~\cite{CT}. It retains the complete source, every
affected neighboring root, the original compact denominators, the full
Casimir row, and all source grades generated by the compact integral. It
does not estimate the variance of the retained source under its marginal
law. It does not identify the sharp rare/good observation in the
companion notes with this observation, or with R-II. It does not
prove mixed-domain admission or S94.

The main new conclusion is a precise field-quadratic-jet decomposition
at the flat retained bath:

\[
 D_y^2 R_{X,f}(1)
 =\beta E_X^* K_{D,f} E_X+\mathcal E_{\beta,f},
 \qquad
 |\mathcal E_{\beta,f}[y,y]|
 \le C b_f\langle y,S_Xy\rangle.                 \tag{1}
\]

Here \(R_{X,f}=\mathbb E_{\mu_\beta}(R_f\mid X)\) is the literal
conditional mean, \(K_{D,f}=D^2D[f](1)\) is the original
directional-square quadratic form, \(E_X\) is the harmonic extension
into the separated eliminated links for the original flat Wilson Hessian
\(H_0=D^2S_0(1)\), and \(S_X=H_0/H_{0,EE}\) is its Schur complement. The
source constant and the full Casimir term are not removed: the latter
and the compact fluctuation correction are inside
\(\mathcal E_{\beta,f}\), and the scalar value of the complete
generating integral remains exact.

The coefficient in the error is \(O(1)\), rather than \(O(\beta)\),
uniformly for \(\beta\ge1\). Moreover,

\[
 |K_{D,f}[a,a]|\le Cb_f\|H_0a\|^2,
 \qquad H_0E_Xy=(0,S_Xy),
\]

and hence

\[
 \boxed{
 |D^2R_{X,f}(1)[y,y]|
 \le Cb_f\{\beta\|S_Xy\|^2+\langle y,S_Xy\rangle\}.}
                                                        \tag{2}
\]

This identifies the low-derivative mixing scale in the true compact
conditional mean: the leading classical term has two more lattice
derivatives than the curvature energy, while the compact/Casimir
remainder has the curvature-energy order with no factor \(\beta\). The
theorem holds for every fixed \(SU(N)\). An exact \(SU(2)\) calculation
below provides a closed-form check.

Only real bath derivatives \textbf{at the flat bath} are asserted. The
source parameters may lie in a fixed complex disc/polydisc. This is not
a complex bath tube, a global zero-free partition theorem, or an
inequality for the full nonlinear marginal distribution.

\hypertarget{original-source-and-actual-observation}{%
\subsection{Original source and actual
observation}\label{original-source-and-actual-observation}}

Use the original Wilson action

\[
 S_0(U)=\sum_p\{N-\operatorname{Re}\operatorname{Tr}U_p\},
 \qquad d\mu_\beta=Z_\beta^{-1}e^{-\beta S_0}\prod_e dU_e,
 \quad \beta\ge1.
\]

The group and its bi-invariant metric are fixed. Choose a real
orthonormal basis \(X_A\) of \(\mathfrak{su}(N)\), and write

\[
 -\operatorname{Re}\operatorname{Tr}(X_AX_B)=\kappa\delta_{AB},
 \quad \sum_A X_A^2=-C_FI,
 \quad \kappa(N^2-1)=C_FN.                       \tag{3}
\]

The literal original ordinary mark is

\[
 R_f=\sum_c f_c r_c^0
 =\beta D[f]-\ell[f]/3,
 \qquad b_f=\sup_c|f_c|.
                                                        \tag{4}
\]

Each \(d_c\) in \(D[f]=\sum_cf_cd_c\) is the original complete
color-summed directional-force-square atom. The atom \(\ell_c\) is its
full Casimir partner. Original pointed offsets and neighboring roots
stay in the determining records. The weights are deterministic scalars
and do not vary within a Lie-algebra color sum.

Let \(E\) be a deterministic separated set of original links,
sufficiently separated that (a) no original plaquette reads two links of
\(E\), and (b) no complete source determining set reads two links of
\(E\). Let \(X\) record every link outside \(E\). This is precisely an
allowed instance of the deterministic separated-link
observation of~\cite{CT}. It is not the observation retaining a random sharp good/bad
geometry.

All supports below have a fixed size in original lattice units. Assume
the torus is larger than the fixed determining collars so they embed in
contractible cubical boxes. Smaller tori are outside this theorem; no
large-torus estimate is obtained by shrinking a collar through a
periodic identification.

Partition the source atoms according to their literal readsets. Atoms
reading no eliminated link belong to \(R_{\mathrm{ext},f}(y)\). Every
other atom belongs to the unique \(R_{e,f}(u_e,y)\) for the eliminated
link it reads. Thus

\[
 R_f=R_{\mathrm{ext},f}(y)+\sum_{e\in E}R_{e,f}(u_e,y).
                                                        \tag{5}
\]

This includes all affected neighboring roots, not only the root at the
eliminated link. The number of atoms and retained coordinates in each
summand is bounded by the actual fixed source stencil.

The conditional Wilson action also splits exactly:

\[
 S_0=S_{\mathrm{ext}}(y)+\sum_{e\in E}S_e(u_e,y).
                                                        \tag{6}
\]

For marks \(f_i\) and source parameters \(\eta_i\), define the
\textbf{raw} source-dependent marginal numerator

\[
 \mathcal W_X(\eta;y)
 =e^{-\beta S_{\mathrm{ext}}(y)+\sum_i\eta_iR_{\mathrm{ext},f_i}(y)}
 \prod_{e\in E} Z_e(\eta;y),
\]

\[
 Z_e(\eta;y)
 =\int_G e^{-\beta S_e(u,y)+\sum_i\eta_iR_{e,f_i}(u,y)}\,du.
                                                        \tag{7}
\]

No \(Z_e\) is dropped or replaced. The fixed original factor
\(Z_\beta^{-1}\) can be restored outside (7); it is independent of \(y\)
and \(\eta\). A raw parent mass or an independent saved scalar remains
outside exactly once. In particular,

\[
 [\eta_i]\log\mathcal W_X(\eta;y)=R_{X,f_i}(y),
                                                        \tag{8}
\]

and higher coefficients are the actual conditional cumulants of these
literal ordinary marks, plus the original source-linear exterior
contribution. This generating function is used to analyze the ordinary
row in estimate (11). It is not asserted to equal the entire physical
source kernel after current conjugation: any original primitive
higher-grade contact, current, reference, or boundary partners remain
additional identified rows. No such partner is deleted or replaced by an
ordinary cumulant.

\hypertarget{flat-one-link-structure-for-every-fixed-sun}{%
\subsection{Flat one-link structure for every fixed
SU(N)}\label{flat-one-link-structure-for-every-fixed-sun}}

At the flat retained bath \(y=1\), each of the six incident plaquettes
reads \(u\) or \(u^{-1}\), and

\[
 S_e(u,1)=6\{N-\operatorname{Re}\operatorname{Tr}u\}=:6s(u).
                                                        \tag{9}
\]

Put \(\varphi_A(u)=\operatorname{Re}\operatorname{Tr}(X_Au)\) and
\(\mathcal Q(u)=\sum_A\varphi_A(u)^2\). Every original directional-force
component affected by this one free link is an integer multiple, with
its actual orientation sign, of \(\varphi_A(u)\). This follows directly
by differentiating a plaquette trace and cyclically moving the generator
next to the single free link. Also

\[
 \operatorname{Re}\operatorname{Tr}(X_Au^{-1})
 =-\operatorname{Re}\operatorname{Tr}(X_Au).
\]

Consequently, after assigning all complete affected atoms,

\[
 D_{e,f}(u,1)=a_{e,f}\mathcal Q(u),
 \qquad
 \ell_{e,f}(u,1)=b_{e,f}\operatorname{Re}\operatorname{Tr}u+c_{e,f},
                                                        \tag{10}
\]

with \(|a_{e,f}|+|b_{e,f}|+|c_{e,f}|\le Cb_f\). These coefficients
include the literal root count and source offsets. They are not
universal coefficients obtained by keeping only the central root.

Both \(s\) and \(\mathcal Q\) are conjugation invariant and invariant
under \(u\mapsto u^{-1}\). Since \(\operatorname{Tr}X_A=0\),

\[
 \mathcal Q(u)\le C\|u-I\|_{\mathrm{HS}}^2=2Cs(u).
                                                        \tag{11}
\]

The complete centered local source therefore satisfies

\[
 |R_{e,f}(u,1)-R_{e,f}(I,1)|\le Cb_f\beta s(u).
                                                        \tag{12}
\]

The value \(R_{e,f}(I,1)\), including its Casimir constant, is kept as
the factor \(\exp\{\sum_i\eta_iR_{e,f_i}(I,1)\}\) in the raw integral.
It is factored only for the following estimates and restored exactly.

\hypertarget{uniform-complex-source-disc-and-compact-moments}{%
\subsection{Uniform complex source disc and compact
moments}\label{uniform-complex-source-disc-and-compact-moments}}

There is \(\delta>0\), depending only on the fixed group and the
original bounded stencil, such that all local normalizers in (7) are
nonzero at \(y=1\) whenever

\[
 \sigma:=\sum_i|\eta_i|b_{f_i}<\delta.             \tag{13}
\]

This radius is independent of \(\beta\ge1\), volume, and the number of
eliminated links.

To prove this, use the actual normalized flat one-link law

\[
 dK_0(u)=\frac{e^{-6\beta s(u)}du}{\int_Ge^{-6\beta s(v)}dv}.
\]

For any fixed \(c<6\), compactness and the nondegenerate unique minimum
\(s(I)=0\) give

\[
 K_0[e^{c\beta s}]\le C_c,
 \quad K_0[(\beta s)^k e^{c\beta s}]\le C_{c,k}.     \tag{14}
\]

Here is an elementary proof of the scale: in a fixed exponential chart
at \(I\), \(s(e^\xi)\) is bounded above and below by fixed multiples of
\(|\xi|^2\), and the Haar Jacobian is bounded above and below. The
denominator is at least \(c\beta^{-d_G/2}\). Split the numerator into
the chart, where Gaussian radial integration applies, and its
complement, where \(s\) has a positive minimum and an exponentially
small bound is stronger than any required power. This also covers
\(1\le\beta\le\beta_0\) after enlarging fixed constants.

Let \(G_\eta=\sum_i\eta_i\{R_{e,f_i}(u,1)-R_{e,f_i}(I,1)\}\). Equations
(12)--(14) show

\[
 |K_0e^{G_\eta}-1|
 \le K_0(|G_\eta|e^{|G_\eta|})\le C\sigma.
\]

Choose \(\delta\) so this is at most \(1/2\). This proves nonvanishing
after restoring the exact scalar factor. It is a local conditional
statement: integrating the retained bath can still create zeros, and no
full-volume partition conclusion follows.

The normalized complex conditional integrals obey the corresponding
total-variation moment bounds. In a symmetric exponential chart, with an
inversion-invariant cutoff,

\[
 |\mathbb E_\eta[\psi(\xi)]|
 \le C\beta^{-k/2}
 \quad\text{if }|\psi(\xi)|\le C|\xi|^k,
                                                        \tag{15}
\]

up to exponentially small exterior terms. This notation is an analytic
normalized integral; positivity or Jensen is not used for complex
\(\eta\). The density and Haar Jacobian are even in \(\xi\), so its odd
local moments vanish exactly with the symmetric cutoff.

Write the one-link exponent as

\[
 V_{e,\eta}(u,y)=\beta L_{e,\eta}(u,y)+J_{e,\eta}(u,y),
\]

\[
 L_{e,\eta}=S_e-\sum_i\eta_iD_{e,f_i},
 \qquad J_{e,\eta}=\frac13\sum_i\eta_i\ell_{e,f_i}.
                                                        \tag{16}
\]

At \((u,y)=(I,1)\), let \(A_{e,\eta}=D_u^2L_{e,\eta}(I,1)\). It is
invertible on the free link algebra and uniformly comparable to the
fixed positive scalar matrix \(6\kappa I\) when (13) holds. The
comparison for complex parameters means an invertible small perturbation
with positive real part on real vectors.

The compact moment bound improves to

\[
 \mathbb E_\eta[\xi\xi^T]
 =\beta^{-1}A_{e,\eta}^{-1}+O(\beta^{-2}).           \tag{17}
\]

For completeness, this does not need an unproved multidimensional
asymptotic theorem. In the symmetric chart, \(L_{e,\eta}\) and
\(J_{e,\eta}\) are even. Thus \(\nabla L=A_{e,\eta}\xi+O(|\xi|^3)\),
\(\nabla J=O(|\xi|)\), and \(\nabla\log J_{\mathrm{Haar}}=O(|\xi|)\).
Integrate the derivative of \(\chi(\xi)\xi_i e^{-V}J_{\mathrm{Haar}}\)
in \(\xi_j\). The derivative of the cutoff is supported in a region with
exponentially small actual mass. Equation (15) bounds the cubic
remainder by \(\beta\mathbb E_\eta|\xi|^4=O(\beta^{-1})\), and the lower
and Haar terms by \(O(\beta^{-1})\). Hence

\[
 \beta\mathbb E_\eta[\xi_i(A_{e,\eta}\xi)_j]
 =\delta_{ij}+O(\beta^{-1}),
\]

which proves (17). All statements are uniform and analytic on any
smaller fixed disc in (13).

\hypertarget{exact-differentiated-normalizer-and-its-schur-term}{%
\subsection{Exact differentiated normalizer and its Schur
term}\label{exact-differentiated-normalizer-and-its-schur-term}}

Let \(U_{e,\eta}(y)=-\log Z_e(\eta;y)\), taking the branch continued
from source zero at the flat bath. For any actual retained coordinates
\(y_a,y_b\), differentiation of the complete compact integral gives

\[
 \partial_aU_{e,\eta}=\mathbb E_\eta(V_a),
 \qquad
 \partial_{ab}U_{e,\eta}
 =\mathbb E_\eta(V_{ab})
  -\operatorname{Cov}_\eta(V_a,V_b).              \tag{18}
\]

The covariance is bilinear for complex sources. These are the actual
induced scores and normalizer responses. They are not replaced by the
original Wilson score.

At the flat bath, expand the smooth finite-dimensional original
functions in the free-link coordinate \(\xi\):

\[
 L_a(\xi,1)=B_a\xi+Q_a(\xi,\xi)+O(|\xi|^3),
 \quad B_a=D_aD_uL(I,1),
\]

\[
 L_{ab}(\xi,1)=C_{ab}+\Lambda_{ab}\xi+O(|\xi|^2),
 \quad C_{ab}=D_aD_bL(I,1).
                                                        \tag{19}
\]

There is no constant term in \(L_a\): at the completely flat field, both
the action and every original directional square have vanishing first
derivatives. The coefficients and the required remainder derivatives are
uniformly bounded on the fixed source disc and fixed determining collar.
The same statements apply to \(J\), with no factor \(\beta\).

Odd-moment cancellation and (15)--(17) yield

\[
 \mathbb E_\eta L_{ab}=C_{ab}+O(\beta^{-1}),
\]

\[
 \operatorname{Cov}_\eta(L_a,L_b)
 =\beta^{-1}B_aA_{e,\eta}^{-1}B_b^T+O(\beta^{-2}).
                                                        \tag{20}
\]

In the second relation, the mixed linear/quadratic term is odd; the
remaining errors are controlled by fourth moments. Subtracting the means
only changes the \(O(\beta^{-2})\) remainder. Mixed \(L,J\) covariances
are \(O(\beta^{-1})\), and \(J,J\) covariances are bounded. Substituting
in (18) proves

\[
 D_y^2U_{e,\eta}(1)
 =\beta\{C-B^TA_{e,\eta}^{-1}B\}
   +\mathcal R_{e,\beta}(\eta),
 \qquad \|\mathcal R_{e,\beta}(\eta)\|\le C.       \tag{21}
\]

The leading bracket is the actual Schur complement of the classical
quadratic action \(S_e-\sum_i\eta_iD_{e,f_i}\). Every full Casimir
contribution is in \(J\), and consequently in the remainder; it has not
been replaced by its value at the minimum.

The scalar \(U_{e,\eta}(1)\), including the exact compact
determinant-type response, is \textbf{not} replaced by a Gaussian scalar
or suppressed. Equation (21) is a statement about the field Hessian of
the exact logarithm.

\hypertarget{uniform-volume-control-of-the-local-remainder}{%
\subsection{Uniform-volume control of the local
remainder}\label{uniform-volume-control-of-the-local-remainder}}

The finite-matrix bound in (21) alone is insufficient for low-momentum
control. Gauge invariance supplies the missing derivative factor.

Each local compact integral is gauge invariant under the original vertex
transformations of all links in its record. Haar invariance of the
free-link integral proves this without fixing a gauge. Its first field
derivative at the flat bath vanishes because a global conjugation would
otherwise give a nonzero invariant linear functional on the simple Lie
algebra \(\mathfrak{su}(N)\).

Consequently, its field Hessian kills infinitesimal gauge gradients,
including mixed pairings with arbitrary tangent fields. The same is true
of the corresponding classical local Schur form. Thus every remainder
\(\mathcal R_{e,\beta}(\eta)\), extended to a full local original-link
box by ignoring the eliminated coordinate, annihilates all original
gradients.

Choose a fixed contractible cubical box containing the complete
determining collar. A finite-box gauge Poincare inequality is

\[
 \inf_\omega\|a-d\omega\|_{\mathrm{collar}}^2
 \le C\sum_{p\text{ in box}}|(da)_p|^2.           \tag{22}
\]

It can be proved by fixing a spanning-tree gauge and expressing each
remaining edge variable through the finitely many plaquette curls in the
box. Its constant depends only on the fixed box, not on \(\beta\) or
volume. Combining (21), gauge annihilation, and (22) gives

\[
 |\mathcal R_{e,\beta}(\eta)[a,a]|
 \le C\sum_{p\text{ in box}_e}|(da)_p|^2.        \tag{23}
\]

Equivalently, on this finite box the bilinear form factors as
\(d^*A_{e,\beta}(\eta)d\), with \(\|A_{e,\beta}(\eta)\|\le C\). This
factorization uses the \textbf{original full box plaquettes}, not merely
retained elementary plaquettes: eliminating a link can leave longer
cycles. Locally flat harmonic/toron directions also vanish because their
restrictions to the contractible box are gradients.

Let

\[
 H_\eta=D^2\left(S_0-\sum_i\eta_iD[f_i]\right)(1),
 \qquad S_{X,\eta}=H_\eta/H_{\eta,EE}.
                                                        \tag{24}
\]

The block \(H_{\eta,EE}\) remains block diagonal and invertible:
separation applies to every original source factor as well as every
action factor. Define \(H_0=H_{\eta=0}\), \(S_X=S_{X,0}\), and

\[
 E_Xy=(-H_{0,EE}^{-1}H_{0,EX}y,\ y).
                                                        \tag{25}
\]

Then

\[
 H_0E_Xy=(0,S_Xy),
 \qquad
 \langle E_Xy,H_0E_Xy\rangle=\langle y,S_Xy\rangle.
                                                        \tag{26}
\]

Sum (21), add the untouched exterior action and source, use bounded
overlap of the actual fixed collars, and evaluate (23) on \(a=E_Xy\).
The result is

\[
 \boxed{
 D_y^2[-\log\mathcal W_X(\eta;\cdot)](1)
 =\beta S_{X,\eta}+\mathcal R_{X,\beta}(\eta),
 \quad
 |\mathcal R_{X,\beta}(\eta)[y,y]|
 \le C\langle y,S_Xy\rangle.}
                                                        \tag{27}
\]

This is an actual compact-law formula on the fixed source disc (13).
Gauge and harmonic zero modes remain zero; no inverse is taken on those
modes. The only inverted block is the uniformly positive fixed-size
free-link block.

\hypertarget{first-ordinary-grade-and-higher-ordinary-exponential-grades}{%
\subsection{First ordinary grade and higher ordinary exponential
grades}\label{first-ordinary-grade-and-higher-ordinary-exponential-grades}}

For one source \(\eta R_f\), differentiation of (24) at zero gives the
exact Schur derivative

\[
 \left.\partial_\eta S_{X,\eta}\right|_0
 =-E_X^*K_{D,f}E_X,
 \qquad K_{D,f}=D^2D[f](1).                       \tag{28}
\]

This follows either by differentiating the block inverse, or by
differentiating the minimized quadratic energy; the stationary free
coordinate eliminates the derivative of the minimizer.

Because \(\partial_\eta[-\log\mathcal W_X]|_0=-R_{X,f}\), (27)--(28)
prove (1). Cauchy's inequality on a slightly smaller source disc gives
the \(Cb_f\) bound for \(\partial_\eta\mathcal R_{X,\beta}\). This
differentiates the actual compact denominator and keeps its
score/contact response.

There is a per-original-root form of (1):

\[
 D^2\mathbb E(r_c^0\mid X)(1)
 =\beta E_X^*D^2d_c(1)E_X+\mathcal E_{\beta,c},
 \qquad \mathcal E_{\beta,c}=E_X^*d^*A_{\beta,c}dE_X,
 \quad \|A_{\beta,c}\|\le C,                       \tag{29}
\]

with a fixed actual collar. This is obtained by activating only that
original source atom before differentiation. It retains every
neighboring-root contribution when the atoms are recombined; it is not a
central-root substitution.

For a grade with \(r\) labeled ordinary marks in (7), Cauchy estimates
give

\[
 \left|
 D_{\eta_1}\cdots D_{\eta_r}
 \mathcal R_{X,\beta}(0)[y,y]
 \right|
 \le C^{r+1}r!\left(\prod_{i=1}^r b_{f_i}\right)
       \langle y,S_Xy\rangle.                    \tag{30}
\]

For distinct squarefree labels the derivative equals the corresponding
coefficient. Taking radii \(\delta/(r b_{f_i})\) gives the factor
\(r^r\), which is bounded by \(e^r r!\); repeated labels have the usual
single-variable factorial. The leading grades are the exact derivatives
of \(\beta S_{X,\eta}\). All source-independent and source-dependent
scalar normalizers remain in (7). Equation (30) concerns the field
Hessians at the flat bath and does not claim a complex bath admission
radius, or identify an arbitrary inherited physical/reference grade with
these ordinary conditional cumulants.

\hypertarget{extra-low-spectral-suppression-of-the-ordinary-quadratic-jet}{%
\subsection{Extra low-spectral suppression of the ordinary quadratic
jet}\label{extra-low-spectral-suppression-of-the-ordinary-quadratic-jet}}

At the flat original field, the directional forces are the directional
divergences of the linearized curvature. In dimensionless Fourier
notation, let

\[
 q_\mu=e^{ip_\mu}-1,
 \quad s(p)=\sum_\mu|q_\mu|^2,
 \quad F_{\mu\nu}=q_\mu a_\nu-q_\nu a_\mu.
\]

With fixed group/metric factors suppressed only in the comparison
constant,

\[
\sum_{\nu,\mu\ne\nu}|q_\mu|^2|F_{\mu\nu}|^2
 =\sum_{\mu<\nu}(|q_\mu|^2+|q_\nu|^2)|F_{\mu\nu}|^2
 \le s(p)\sum_{\mu<\nu}|F_{\mu\nu}|^2
 =\|(sI-qq^*)a\|^2.                            \tag{31}
\]

Thus \(K_{D,1}\preceq CH_0^2\), and signed deterministic marks obey
\(|K_{D,f}[a,a]|\le Cb_f\|H_0a\|^2\). No directional-score identity is
used: this is a flat quadratic Hodge relation. Equations (1) and (26)
yield (2).

On the spectral subspace \(0<S_X\le\varepsilon\),

\[
 |D^2R_{X,f}(1)[y,y]|
 \le Cb_f(\beta\varepsilon+1)\langle y,S_Xy\rangle.
                                                        \tag{32}
\]

Relative to the leading marginal action \(\beta S_X\), the scalar source
response is of order \(b_f(\varepsilon+\beta^{-1})\). The \(\beta^{-1}\)
contribution is the retained compact/Casimir mixing term. Neither the
shared retained means nor a projected cross term has been removed by
this estimate.

The spectral cutoff here is for the deterministic flat \textbf{spatial}
operator \(S_X\). It is not a low-spectrum assertion for the original
heat-bath generator or for the arbitrary inherited output projection
\(Q\) in estimate (11).

\hypertarget{exact-source-zero-induced-hessian-su2-check}{%
\subsection{Exact source-zero induced Hessian; SU2
check}\label{exact-source-zero-induced-hessian-su2-check}}

At the flat bath define

\[
 dK_\beta(u)=z_\beta^{-1}e^{6\beta\operatorname{Re}\operatorname{Tr}u}\,du,
 \qquad m_N(\beta)=N^{-1}K_\beta\operatorname{Re}\operatorname{Tr}u.
\]

Conjugation and inversion symmetry give \(K_\beta u=m_NI\). Monotonicity
of the mean trace in \(\beta\), its zero Haar mean at \(\beta=0\), and
nonconstant trace give

\[
 0<m_N(1)\le m_N(\beta)\le1,
 \qquad \beta(1-m_N(\beta))\le C_N.               \tag{33}
\]

The last bound follows from (14). Compact integration by parts gives
exactly

\[
 K_\beta[\varphi_A\varphi_B]
 =\frac{\kappa m_N(\beta)}{6\beta}\delta_{AB}.
                                                        \tag{34}
\]

Indeed \(L_A\operatorname{Re}\operatorname{Tr}u=\varphi_A\) and
\(K_\beta L_A\varphi_B=-\kappa m_N\delta_{AB}\), while the logarithmic
density derivative is \(6\beta\varphi_A\).

For six staple velocities \(x_i\) at the flat bath, the exact star
contribution to the induced action Hessian is therefore

\[
 D^2U_e(1)[x,x]
 =\beta\kappa m_N\left\{
     \sum_{i=1}^6|x_i|^2-\frac16\left|\sum_{i=1}^6x_i\right|^2
   \right\}.                                          \tag{35}
\]

This is \(m_N\) times the classical harmonic star Schur complement. It
holds after any actual exterior-link parameterization: the normal second
derivative of each group staple supplies the displayed quadratic term,
while tangent accelerations pair with a scalar mean matrix and have zero
trace.

The raw star action constant is \(6N\), independent of the bath, so its
factor \(e^{-6N\beta}\) is retained and has zero bath Hessian. The
retained exterior action is still present. For the entire separated set,

\[
 D^2[-\log\mathcal W_X(0;\cdot)](1)
 =\beta H_{\mathrm{ext}}+
   \beta m_N\sum_{e\in E}H_{e,\mathrm{Schur}},
\]

\[
 m_N\beta S_X\preceq D^2[-\log\mathcal W_X(0;\cdot)](1)
 \preceq\beta S_X.                                   \tag{36}
\]

For \(SU(2)\), write \(u\in S^3\), six staples \(v_i\),
\(H=\sum_i v_i\). Then

\[
 Z_{e,\mathrm{raw}}=e^{-12\beta}
 \frac{2I_1(2\beta|H|)}{2\beta|H|}.
\]

At \(v_i=e_0\), \(t=12\beta\), \(m=I_2(t)/I_1(t)\), and \(q=m/t\). In
Euclidean quaternion tangent coordinates, \(\kappa=2\), so (35) is

\[
 2\beta m\sum_i|x_i-\bar x|^2.                       \tag{37}
\]

One has \(m\ge3/4\) for \(\beta\ge1\): the Bessel recurrence gives
\(m=t/(4+tI_3(t)/I_2(t))\ge t/(t+4)\), since \(I_3/I_2<1\) follows from
its positive spherical-mean representation.

For the literal central rooted ordinary member, the finite-volume exact
formula gives

\[
 D^2\Phi_e(1)[x,x]
 =2c_X^2\beta(1-q)
 \sum_{\mu=1}^3\left|
 x_{\mu+}+x_{\mu-}-\frac13\sum_i x_i
 \right|^2.
                                                        \tag{38}
\]

Since the sum of these pair-average differences is bounded by
\(2\sum_i|x_i-\bar x|^2\),

\[
 D^2\Phi_e(1)\preceq
 \frac{2c_X^2(1-q)}m\,D^2U_e(1).                 \tag{39}
\]

Equation (38) is a check for the rooted member only. The complete-source
results are (1), (27), and (29); they include the neighboring roots.

There is also a useful exact complete-readout formula on an arbitrary
real \(SU(2)\) bath. Every affected original ordinary atom is at most
quadratic in the one free quaternion, since each original plaquette and
invariant first derivative is linear in that quaternion. Thus the
complete assigned source has the form

\[
 R_{e,f}(u,y)=A(y)+b(y)\cdot u+u^TC(y)u,
 \qquad C=C^T.
\]

For \(H=rn\), put \(m=I_2(2\beta r)/I_1(2\beta r)\) and
\(q=m/(2\beta r)\). Its \textbf{exact full conditional mean} is

\[
 A+m\,b\cdot n+q\operatorname{tr}C+(1-4q)n^TCn.   \tag{39a}
\]

This formula includes all original affected roots through \(A,b,C\). At
\(H=0\) use the Haar limits \(m=0,q=1/4\); the expression is independent
of the unused direction \(n\). Higher source grades are obtained by
integrating the corresponding polynomials of degree at most \(2r\)
against the same raw compact integral, or equivalently by its
derivatives in \(H\); its denominator is retained.

As a literal one-link check for the uniform original source, when only
one link differs from the identity, the complete directional-square sum
has coefficient \(30c_X^2(1-u_0^2)\): the central root contributes
\(12c_X^2\), and the eighteen other incident roots contribute
\(18c_X^2\). This agrees with the finite-volume complete two-link quadratic
coefficient. The entire Casimir partner changes by \(-36c_X^2(1-u_0)\)
in \(\ell\), so the ordinary source variation is

\[
 R(U)-R(I)=30c_X^2\beta(1-u_0^2)+12c_X^2(1-u_0).
                                                        \tag{40}
\]

The exact constants \(R(I)\), the untouched root terms, and the scalar
compact integral remain in the raw generating expression.

\hypertarget{extension-to-real-flat-backgrounds-and-their-torus-holonomies}{%
\subsubsection{Extension to real flat backgrounds and their torus
holonomies}\label{extension-to-real-flat-backgrounds-and-their-torus-holonomies}}

The flat-bath assertions extend to any real original lattice connection
\(A\) whose elementary plaquette holonomies are all the identity. This
includes arbitrary commuting torus holonomies. Every fixed contractible
determining collar has a genuine vertex gauge transformation taking that
flat connection to the identity on the collar. Apply that transformation
also to the integrated free link and use exact Haar invariance. It
follows that each local raw integral \(Z_e(\eta;A_X)\), its complete
ordinary source grades, and the exterior source constants have the same
scalar values as at the identity bath. Therefore

\[
 \mathcal W_X(\eta;A_X)=\mathcal W_X(\eta;1)
                                                        \tag{FH1}
\]

for the separated compact integration and ordinary marks in (7), with
the identical original scalar factors. In particular, \(R_{X,f}(A_X)\)
and all higher ordinary conditional cumulants at these exactly flat
baths are independent of the flat holonomy. An arbitrary mixture of such
flat bath values creates no between-flat-bath variance of these local
conditional means.

The field-jet formulas hold with the original covariant flat Hessian
\(H_A\), its separated-link harmonic extension \(E_{X,A}\), and its
Schur complement \(S_{X,A}\). The gauge transformations on collars are
unitary on tangent colors and preserve every local constant. Flat
covariant differences commute, so the Fourier/Hodge argument in (31)
gives

\[
 |K_{D,f,A}[a,a]|\le Cb_f\|H_Aa\|^2
\]

with the same constant, retaining all covariant gauge and harmonic zero
modes. Thus (1), (2), and (27) hold uniformly over these flat
backgrounds.

Equation (FH1) concerns \textbf{actual compact integration of the
separated free links at an exactly flat retained bath}. It does not
identify that conditional law with a Gaussian field fluctuating over the
whole torus. A full Gaussian reference can still have finite-torus
holonomy corrections to its fluctuation means, as in the finite-volume
benchmark. Nor does (FH1) bound the distribution or shared means of
actual nonflat histories.

\hypertarget{uniform-real-chart-third-derivatives-of-the-actual-compact-mean}{%
\subsection{Uniform real-chart third derivatives of the actual compact
mean}\label{uniform-real-chart-third-derivatives-of-the-actual-compact-mean}}

The preceding theorem concerns flat jets. A separate local regularity
estimate is available on a fixed \textbf{real} near-flat bath chart and
can be used for an actual-law local Taylor remainder. No complex bath
admission or sharp selected-domain statement is made here.

For each complete original atom \(r_c^0\), let
\(m_c(y)=\mathbb E_{\mu_\beta}(r_c^0\mid X=y)\). There are fixed
bath-chart radii \(0<\delta_1<\delta_0\), depending only on the group
and the literal determining collar, such that

\[
 \sup_{|y|<\delta_1}\|D_y^j m_c(y)\|\le C_j\beta,
 \qquad j=0,1,2,3,\quad\beta\ge1.                \tag{LC1}
\]

These are direct original-link coordinates in the fixed local bath
collar, or a fixed gauge transform thereof; the transformed prepared
readout is differentiated completely. At the flat point
\(D_y m_c(0)=0\), by the gauge/global-adjoint argument already given. An
atom reading no eliminated link has (LC1) directly from its original
smooth formula. Otherwise it reads one free link, and the following
argument applies.

At the identity bath, the one-link action \(S_e(u,0)=6s(u)\) has a
unique global minimum at \(u=I\) with positive Hessian. Compactness
gives a strictly positive action separation outside a fixed identity
neighborhood. For all real \(|y|<\delta_0\), the implicit function
theorem and this action separation give a unique global minimizer
\(u_*(y)\), analytic with uniformly bounded derivatives on a smaller
chart. The local free-link Hessian remains uniformly positive. This is a
fixed-dimensional one-link fact; no global gauge-field convexity is
used.

Make the exact Haar change of variables \(u=u_*(y)v\), and put

\[
 L(v,y)=S_e(u_*(y)v,y)-S_e(u_*(y),y).
\]

The raw factor \(\exp\{-\beta S_e(u_*(y),y)\}\) remains in the original
denominator. It cancels only in this normalized conditional mean. In a
fixed symmetric logarithm chart \(v=e^\xi\),

\[
 L(0,y)=0,\qquad D_\xi L(0,y)=0
\]

for every real bath in the chart. Thus, for \(1\le j\le3\),

\[
 |D_y^jL(\xi,y)|\le C_j|\xi|^2,
 \qquad L(\xi,y)\ge c|\xi|^2.                    \tag{LC2}
\]

Outside a fixed smaller free-link chart there is a uniform positive
action separation. The exact normalized kernel consequently has

\[
 \|\beta D_y^jL\|_{L^p(K_y)}\le C_{j,p}
 \quad (j\le3, p<\infty),                      \tag{LC3}
\]

uniformly in \(y\) and \(\beta\ge1\). Inside the chart this follows from
the Gaussian-scale moment bound. Outside, any fixed power of \(\beta\)
is dominated by the exponentially small mass, with the true compact
normalizer bounded below by \(c\beta^{-d_G/2}\) after the minimum is
factored.

The complete transported carrier \(A_c(v,y)=r_c^0(u_*(y)v,y)\) satisfies
\(\sup_v|D_y^jA_c(v,y)|\le C_j\beta\) for \(j\le3\). This includes every
derivative of \(u_*(y)\) and the original readout. Differentiating the
exact normalized integral uses

\[
 D_a K_y(A)=K_y(D_aA)-\operatorname{Cov}_{K_y}(A,\beta D_aL).
                                                        \tag{LC4}
\]

Repeated differentiation through order three gives finitely many moments
of derivatives of \(A_c\) multiplied by products of the score
derivatives \(\beta D_y^jL\), with the exact differentiated denominator
terms. Equations (LC3) and Hölder bound each term by \(C\beta\), proving
(LC1). There is no moving integration boundary after the Haar
translation, and no density has been reset to Wilson at a later layer.

In particular, on this real bath chart,

\[
 \left|m_c(y)-m_c(0)-\frac12D^2m_c(0)[y,y]\right|
 \le C\beta|y|^3.                               \tag{LC5}
\]

To turn (LC5) into an actual original-law moment bound, the bath must
first be placed in such a fixed local chart using an observed-link
spanning tree and the literal original plaquette moment record. A global
gauge choice, native-law comparison, or selector regularity does not
follow from this lemma. This step is addressed separately.

\hypertarget{sharp-selected-next-layers-exact-formulas-and-limits}{%
\subsection{Sharp selected next layers: exact formulas and
limits}\label{sharp-selected-next-layers-exact-formulas-and-limits}}

For the specified sharp observation \(\mathcal B\), condition on a fixed
geometric record \(D\) and write its actual eliminated fiber as
\(\Omega_D(y)\). The raw density after eliminating its rare components
and safe seeds is

\[
w_D(y)=\chi_{\mathrm{ext},D}(y)e^{-\beta S_{\mathrm{ext},D}(y)}
        \prod_C Z_C(D,y)\prod_{s\in S(D)}Z_{s,\mathrm{good}}(D,y),
                                                        \tag{41}
\]

with every surviving exterior selector in \(\chi_{\mathrm{ext},D}\) and
all actual boundary/selector indicators in the integrals defining the
component factors. Further original retained factors, if present, stay
in this same raw expression. When a smooth coordinate differentiation is
legitimate, set

\[
 I_y(G)=\int_{\Omega_D(y)}G(z,y)e^{-V(z,y)}dz.
\]

If \(v_a\cdot n\) is the normal velocity of the moving boundary,
Reynolds' formula gives

\[
 \partial_a I_y(G)=
 \int_{\Omega_D(y)}(\partial_aG-G\partial_aV)e^{-V}dz
 +\int_{\partial\Omega_D(y)}G e^{-V}(v_a\cdot n)d\sigma.
                                                        \tag{42}
\]

With \(m_G=I_y(G)/I_y(1)\), this becomes

\[
 \partial_am_G=
 E_y\partial_aG-\operatorname{Cov}_y(G,\partial_aV)
 +\frac1{I_y(1)}\int_{\partial\Omega_D(y)}
        (G-m_G)e^{-V}(v_a\cdot n)d\sigma.         \tag{43}
\]

Similarly, \(\partial_a\log I_y(1)\) includes the raw boundary mass. A
second derivative includes the differentiated surface terms, moving
intersections/corners, and the induced-score covariance. At a cusp these
formulas are interpreted weakly; a nonexistent classical second
derivative cannot be supplied by the one-link spectral gap.

On a fixed domain, the next marginal potential is
\(V_{\mathrm{eff}}=V_{\mathrm{ext}}-\sum\log Z\), and its Hessian has
the exact bulk contribution

\[
 D^2V_{\mathrm{eff}}
 =D^2V_{\mathrm{ext}}+\sum\{E(D^2V)-\operatorname{Cov}(DV,DV)\},
                                                        \tag{44}
\]

plus the surface terms for moving selected domains. The induced density
is therefore not the original Wilson density. Equations (41)--(44)
retain the existing raw \(Z_C\) and \(Z_{s,\mathrm{good}}\); they give
no independent-color tensorization for a further layer.

Full raw recombination of good and bad sectors restores the smooth
compact integral in (7). It must occur before taking a logarithm or
replacing a conditional mean. For a schematic complete partition
\(Z=\sum_DZ_D\), put \(p_D=Z_D/Z\). Where classical derivatives exist,
the exact identities are

\[
 \partial^2\log Z
 =\sum_Dp_D\partial^2\log Z_D
    +\operatorname{Var}_{p_D}(\partial\log Z_D),
\]

\[
 \partial\Bigl(\sum_Dp_Dm_D\Bigr)
 =\sum_Dp_D\partial m_D
    +\operatorname{Cov}_{p_D}(m_D,\partial\log Z_D). \tag{45}
\]

In a nonsmooth case the identities are understood by first recombining
the raw integrals or through their weak derivative. Between-history
means and scores remain in (45). The preceding flat-jet theorem uses the
recombined compact integral from the outset. It does not apply (27) to
the individual sharp factors in (41), and it does not delete masks after
a later projection.

\hypertarget{a-checked-cusp-for-the-complete-original-ordinary-source}{%
\subsubsection{A checked cusp for the complete original ordinary
source}\label{a-checked-cusp-for-the-complete-original-ordinary-source}}

The distinction persists even at a flat bath, and can be checked without
a rooted-source substitution. In \(SU(2)\), take the six actual staples
to be five copies of \(e_0\) and

\[
 v_1(t)=\cos t\,e_0+\sin t\,e_1.
\]

These staples can be realized by actual distinct exterior links as in
the finite-volume one-link construction. With \(c=1-\epsilon^2/2\in(0,1)\),
the exact all-good star domain is

\[
 \Omega(t)=\{u\in S^3:u_0\ge c,\quad u_0\cos t+u_1\sin t\ge c\}.
\]

Other exterior selectors remain strictly good for sufficiently small
\(|t|\); their original factors are retained. Let \(Z_g(t)\) be the raw
central integral over this domain. A first variation at the cap boundary
gives

\[
 Z_g(t)=Z_g(0)-L|t|+o(|t|),
 \qquad L=\frac{1-c^2}{2\pi}e^{-12\beta(1-c)}>0.    \tag{46}
\]

Indeed the smooth integrand derivative integrates to zero by the
reflection induced by a global conjugation. For \(t>0\), the
intersection loses the boundary strip where \(u_1<0\). Its loss
coefficient is

\[
 \int_{u_0=c,u_1<0}\frac{-u_1}{|\nabla_{S^3}u_0|}
     e^{-12\beta(1-c)}\frac{d\sigma}{2\pi^2}
 =\frac{1-c^2}{2\pi}e^{-12\beta(1-c)}.
\]

At \(t=0\), the \textbf{complete uniform original ordinary sum},
including all neighboring roots, is the function

\[
 F_0(x)=R(I)+30c_X^2\beta(1-x^2)+12c_X^2(1-x),
 \qquad x=u_0.
\]

It is strictly decreasing on \([c,1]\). Thus its genuine selected
conditional mean has

\[
 M_g(t)=M_g(0)
 -\frac{L}{Z_g(0)}\{F_0(c)-M_g(0)\}|t|+o(|t|),
 \qquad F_0(c)>M_g(0).                            \tag{47}
\]

The smooth source and action derivative again integrates to zero by
global conjugation. The original scalar \(R(I)\) is present in both
terms and cancels only in their exact difference. Exterior source terms
vary smoothly and do not remove the displayed cusp. Full raw good/bad
recombination cancels the cap loss in both numerator and denominator.
Equations (46)--(47) prove that the smooth compact flat-jet theorem
cannot be applied to the isolated sharp selected good kernel, even
though that kernel can have a large internal spectral gap. This is an
obstruction to that proof step, not a counterexample to the desired
original-law susceptibility.

\hypertarget{precise-limitation-for-the-retained-variance}{%
\subsection{Precise limitation for the retained
variance}\label{precise-limitation-for-the-retained-variance}}

The results identify a real compact-law low-spectral quadratic mechanism
at the actual target \(R_{X,f}\). They do not prove a bound on

\[
 \left\|Q\,\mathbb E[(R_{X,f}-\mathbb ER_f)\mid Y]\right\|_2^2.
\]

Two direct approaches to a bridge are examined.

\textbf{Outer pressure curvature.} Integrating (7) over the retained
variables is exactly the original source partition function. At source
zero its logarithmic curvature is

\[
 \left.\partial_\eta^2\log\int\mathcal W_X(\eta;y)dy\right|_0
 =\mathbb E\left[\left.\partial_\eta^2\log\mathcal W_X(\eta;X)\right|_0\right]
   +\operatorname{Var}(R_{X,f}).                  \tag{48}
\]

The first term is the conditional internal variance. The separated
compact kernels and available local moment bounds control that extensive
term. The second term is exactly the retained ordinary susceptibility
still required. Nonvanishing of every local \(Z_e\) does not imply
nonvanishing or bounded logarithmic curvature after the retained
integral. Original multipoint exponential energy bounds give uncentered
moment control but do not supply the missing cancellation of the second
term. Also, a source-tilted law is not the Wilson law to which the
moment theorem was originally applied. Consequently this
pressure argument does not close (11).

\textbf{Averaging the quadratic-jet comparison.} Equations (27) and (32)
concern field derivatives at the flat bath. That bath has zero mass
under the retained continuous law. They cannot be inserted as a Hessian
inequality on typical nonlinear configurations. Even a proved averaged
local Hessian error would not by itself control the inverse marginal
covariance/Poisson operator: one would need a coercivity or
source-specific resolvent estimate for that actual law, with its gauge
modes, nonflat centers, selectors and induced normalizers. Neither (14)
nor local zero-freeness proves such an estimate. A gauge-dependent fixed
flat preconditioner also cannot be substituted for a gauge-covariant
nonlinear one without a verified transformation and its responses.

To turn the flat-jet statement into that inequality would require
control of the actual nonlinear marginal over its relevant
configurations and histories, including its induced normalizers and
geometric/boundary responses. A Gaussian quadratic benchmark can now use
the precise \(O(1)\) curvature-order remainder and the fourth-derivative
classical term, but it is still a benchmark until an actual-law
comparison is proved.

The theorem does not set shared history means to zero, change source
grades, replace the prepared readout, or control the difference of
squared projected variances by an error norm. Physical factors
\(c_\beta^{-1}a^{-2}\) and their squares remain outside the
dimensionless estimates with their original labels. No mixed-domain or
absolute-vacuum construction follows from the flat bath calculation
alone.

\newpage

\hypertarget{part-iii.-covariance-bookkeeping-and-the-leading-reference-source}{%
\section{Part III. Covariance bookkeeping and the leading reference
source}\label{part-iii.-covariance-bookkeeping-and-the-leading-reference-source}}

\hypertarget{status-and-dependency-ledger}{%
\subsection{Status and
dependencies}\label{status-and-dependency-ledger}}

This note does \textbf{not} establish the retained ordinary estimate
(11) for the original Wilson law. It proves three more precise
statements:

\begin{enumerate}
\def\labelenumi{\arabic{enumi}.}
\tightlist
\item
  Under the proposed original-law rare-component/safe-seed construction,
  with all action, good/bad selector and complete source determining
  factors included, the covariance removed by the safe good layer is a
  positive semidefinite finite-range operator. The rare removed
  covariance is exponentially localized under the stated component-mass
  and local-moment hypotheses. Consequently the retained and original
  susceptibility operators differ by a uniformly bounded positive
  operator. This is an exact original-law reduction; it supplies no
  missing bound on the retained operator.
\item
  In the complete flat Gaussian benchmark of the
  companion notes, the directional-force covariance decays as
  distance to the power minus six, and the complete ordinary covariance
  as distance to the power minus twelve. The ordinary source has a
  stronger source-specific low-spectral estimate of order epsilon to the
  fourth in four dimensions. Arbitrary mixtures of commuting flat
  holonomies retain an explicitly bounded shared mean term. The
  reference ordinary source conditional on genuine low Fourier modes
  satisfies a corresponding concentration bound with scale epsilon and
  variance factor epsilon\textsuperscript{4+n}-8.
\item
  Even an actual Wilson positive observation can destroy pointwise
  covariance decay of conditional means: conditioning on the total
  source produces a spatially constant conditional-mean covariance. This
  does not contradict an extensive operator-norm estimate. It explains
  why one must retain complete marked sums and cannot transfer rootwise
  covariance estimates through arbitrary histories or later masks.
\end{enumerate}

The Gaussian model is not substituted for the Wilson model. No assertion
identifies the Fourier observation with the observation X of~\cite{CT}, the proposed O/B, or an
R-II integration. The results below concern the original scalar ordinary
row with deterministic color-independent marks and its fixed finite
pointing completion. They do not automatically include arbitrary
inherited source derivatives or saved complex grades.

The old Gaussian total-variance and flat-holonomy mean estimates are
already in the companion notes. They are inputs or reproved
ingredients, not new achievements of this note. The new parts are the
explicit spatial kernel, scale-dependent spectral/concentration
estimates, the exact good-layer covariance operator statement, and the
observation example.

\hypertarget{exact-covariance-subtraction-under-the-proposed-original-law-coarsening}{%
\subsection{Exact covariance subtraction under the proposed original-law
coarsening}\label{exact-covariance-subtraction-under-the-proposed-original-law-coarsening}}

Write the complete dimensionless ordinary source as

\[
 F_f=\sum_c f_c r_c^0,\qquad M_f=\sum_c|f_c|^2,
\]

with every original color incidence, neighboring root, pointed offset
and test transform retained in the local terms r\_c\^{}0. Marks f\_c are
deterministic. All statements in this section are for the original
probability law, with the actual raw component and one-link normalizers
retained in its conditional disintegration.

Let

\[
 \mathcal O=\sigma(D,\text{links outside }E_D),\qquad
 \mathcal B=\sigma(D,\text{links outside }E_D\cup S(D)),
 \qquad \mathcal B\subseteq\mathcal O.
\]

The parentheses in the second formula mean that both E\_D and S(D) are
integrated out. The safe seed set is chosen so that every original
action, good/bad test, and complete source determining factor containing
a seed avoids E\_D. Distinct seeds are separated enough that no such
factor contains two of them. This uses the complete factor graph; a
bad-plaquette adjacency alone is insufficient.

\hypertarget{assumptions-needed-for-the-locality-conclusion}{%
\subsubsection{Assumptions needed for the locality
conclusion}\label{assumptions-needed-for-the-locality-conclusion}}

The following properties are the exact relevant content of the proposed
construction, rather than consequences of notation alone.

\begin{itemize}
\tightlist
\item
  Conditional on O, the eliminated rare components are independent
  compact fibers, including every sharp test and each original raw Z\_C.
\item
  Every original local source root either has no variable in the
  eliminated set or reads one such component. If two eliminated pieces
  occur in one determining factor they have been joined into the same
  component.
\item
  Conditional on B, the same rare component kernels remain, and the safe
  seeds have independent exact sharp good-star kernels with their raw
  Z\_s,good. The safety condition makes every rare component action/test
  normalizer and integrated source mean independent of the free safe
  seeds.
\item
  There is a uniform source determining radius r\_0 and uniformly
  bounded incidence. The original local moment estimate gives
  \(\|r_c^0\|_q\le Cq\) (or the specified fixed source-grade analogue)
  on each fixed determining record.
\item
  The component events obey the proposed bound
  \(p_C\le e^{-\lambda |C|}\), and the number of allowed size-m
  components touching one root is bounded by \(A^m\), with A fixed. The
  enlarged component graph has bounded degree, and a component touching
  roots at distance d has size at least a fixed multiple of d, up to a
  fixed source collar.
\end{itemize}

The note does not independently manufacture these properties from
missing definitions of E\_D. They must hold for the literal construction
before the theorem is applied. In particular, raw normalizers are not
discarded, and good-test couplings are not omitted.

\hypertarget{the-exact-covariance-identities}{%
\subsubsection{The exact covariance
identities}\label{the-exact-covariance-identities}}

Let r denote the vector of centered local roots. Define covariance
matrices as quadratic forms on finitely supported deterministic marks:

\[
 C_{cd}=\operatorname{Cov}(r_c^0,r_d^0),\quad
 C^{\mathcal O}_{cd}=\operatorname{Cov}(E[r_c^0|\mathcal O],E[r_d^0|\mathcal O]),\quad
 C^{\mathcal B}_{cd}=\operatorname{Cov}(E[r_c^0|\mathcal B],E[r_d^0|\mathcal B]).
\]

The law of total covariance and the tower property give exactly

\[
 C=C^{\mathcal O}+K_{\rm rare},\qquad
 C^{\mathcal O}=C^{\mathcal B}+K_{\rm good},
\tag{1.1}
\]

where

\[
 (K_{\rm rare})_{cd}=E\operatorname{Cov}(r_c^0,r_d^0|\mathcal O),
\]

\[
 (K_{\rm good})_{cd}
 =E\operatorname{Cov}(E[r_c^0|\mathcal O],E[r_d^0|\mathcal O]\mid\mathcal B).
\tag{1.2}
\]

Both K matrices are positive semidefinite as operators on marks. Their
individual off-diagonal entries need not be nonnegative.

For the good layer, conditional on B, write the O-conditional root as
its B-measurable part plus its terms depending on individual safe seeds.
The rare-component conditional means belong to the B-measurable part:
their action factors, tests, readouts and normalizers avoid each safe
seed by the safety condition. Distinct seed variables are conditionally
independent. A covariance therefore survives only if both roots read the
same safe seed. Hence

\[
 (K_{\rm good})_{cd}=0
 \quad\text{when }\operatorname{dist}(c,d)>2r_0,
\qquad
 |(K_{\rm good})_{cd}|\le C.
\tag{1.3}
\]

The diagonal bound follows from conditional variance contraction and the
original local L2 moments. Fixed incidence and Schur's test then give

\[
 0\preceq K_{\rm good}\preceq C I.
\tag{1.4}
\]

The random seed geometry does not create a new long-range term in
K\_good: it is conditioned on inside B. Geometry-to-geometry covariance
remains inside C\^{}B, exactly where it belongs.

For the rare part, the same conditional factorization gives
contributions only from components touching both roots. For a particular
component event, conditional Cauchy-Schwarz and the local rare-weighted
moment inequality give

\[
 \left|E[1_{\{C\}}\operatorname{Cov}(r_c^0,r_d^0|\mathcal O)]\right|
 \le C p_C\log^2(e/p_C).
\tag{1.5}
\]

For sufficiently large lambda, the function p log\^{}2(e/p) is
increasing on the relevant interval, so the right side is bounded by

\[
 C(1+\lambda |C|)^2e^{-\lambda|C|}.
\]

Summing the bounded-degree animals gives an exponentially localized
kernel. In particular,

\[
 \sup_c\sum_d |(K_{\rm rare})_{cd}|
 \le C\lambda^2e^{-\lambda},
\tag{1.6}
\]

with a fixed enlarged constant for lambda above a fixed threshold. One
explicit off-diagonal form, after enlarging that threshold, is

\[
 |(K_{\rm rare})_{cd}|
 \le C\lambda^2 e^{-c\lambda\max(1,\operatorname{dist}(c,d)/r_0)}.
\tag{1.7}
\]

The stronger row-sum prefactor in (1.6) follows by grouping first by
component size: the number of roots touched by a size-m component adds
only a polynomial in m, and the m=1 contribution controls the
geometrically convergent sum. The exact constants depend on the literal
enlarged graph and source collars. If its component size convention
changes, the constants in the stated mass bound must be tracked rather
than silently renamed.

Combining the identities gives

\[
 \boxed{\quad
 C=C^{\mathcal B}+K,\qquad
 0\preceq K\preceq C(1+\lambda^2e^{-\lambda})I.
 \quad}
\tag{1.8}
\]

Thus, for every deterministic marked sum,

\[
 0\le\operatorname{Var}(F_f)
       -\operatorname{Var}(E[F_f|\mathcal B])
 \le C(1+\lambda^2e^{-\lambda})M_f.
\tag{1.9}
\]

This is an exact additive variance statement because B is a genuine
conditional expectation of F. It does not use an approximate source
identity. At the operator level,

\[
 \|C^{\mathcal B}\|\le\|C\|
 \le\|C^{\mathcal B}\|+C(1+\lambda^2e^{-\lambda}).
\tag{1.10}
\]

Therefore an unbounded susceptibility eigenvalue cannot be removed by
this single exact layer. Conversely a uniform bound on either of the two
covariance operators supplies one on the other. In a
translation-compatible periodic construction, the removed kernel is a
finite-range part plus an exponentially localized part; its structure
factor is analytic in a fixed momentum strip. Any nonanalytic infrared
part of the retained covariance remains the same as that of the original
covariance, up to that analytic correction.

\hypertarget{later-projections-retain-the-complete-error-not-the-variance-identity}{%
\subsubsection{Later projections retain the complete error, not the
variance
identity}\label{later-projections-retain-the-complete-error-not-the-variance-identity}}

Put e\_f=F\_f-E(F\_f\textbar B). Equation (1.9) is the bound
\(\|e_f\|_2^2\le CM_f\). Any later conditional expectation P\_Y and
orthogonal output projection Q give

\[
 \|QP_Ye_f\|_2^2\le CM_f.
\]

For \(A=QP_Y(E(F_f|B)-EF_f)\) and \(e=QP_Ye_f\), the projected
squared-norm difference is exactly

\[
 \|A+e\|_2^2-\|A\|_2^2
 =2\operatorname{Re}\langle A,e\rangle+\|e\|_2^2.
\tag{1.11}
\]

The original B-orthogonality does not imply that this projected cross
term vanishes. Nor does the complete error bound allow separate later
masks on component pieces. This preserves the restriction stated
in~\cite{CT}.

\hypertarget{a-pointwise-spatial-estimate-for-the-full-gaussian-ordinary-benchmark}{%
\subsection{A pointwise spatial estimate for the full Gaussian ordinary
benchmark}\label{a-pointwise-spatial-estimate-for-the-full-gaussian-ordinary-benchmark}}

Use the fixed-lattice exact flat-holonomy Gaussian model (as in the companion
notes), not the nonlinear Wilson law. At unit lattice spacing and
side n, let V=n\^{}4. In one adjoint color sector of a commuting flat
background, momentum lies on a grid shifted by theta/n.~Put

\[
 p_\mu(k)=e^{ik_\mu}-1,\quad
 u_\mu=|p_\mu|^2,\quad s=\sum_\mu u_\mu,\quad
 K=sI-pp^*,\quad K^+=s^{-1}P_T.
\]

At s=0 the source is zero; exact harmonic modes may be fixed or assigned
any probability law. With the fixed-lattice fixed original normalization
upsilon,

\[
 S_2=\frac\upsilon2\langle A,KA\rangle,
 \qquad \operatorname{Cov}(A|\theta)=(\beta\upsilon)^{-1}K^+.
\]

Up to harmless orientation phases, a directional force row is

\[
 L_{\nu\mu}(k)A
 =\overline p_\mu(p_\mu A_\nu-p_\nu A_\mu),
 \qquad L_{\nu\mu}p=0.
\tag{2.1}
\]

The original linear force is upsilon times this row. The covariance
symbol of the normalized directional forces
\(\sqrt\beta F_{e;\mu}^{(1)}\) is therefore a fixed multiple of

\[
 H_{ij}(k)=L_i(k)K^+(k)L_j(k)^*
          =\frac{L_i(k)L_j(k)^*}{s(k)}.
\tag{2.2}
\]

Its numerator vanishes to fourth order at the origin and its denominator
to second order. Consequently, away from the origin,

\[
 |\partial^J H_{ij}(k)|\le C_J|k|^{2-|J|}.
\tag{2.3}
\]

On a dyadic annulus of radius rho, the L1 norm of its m-th derivative is
bounded by C rho\^{}(6-m). Integrating by parts m\textgreater6 times on
each annulus, and splitting at rho comparable to
(1+\textbar x\textbar)\^{}-1, proves

\[
 |\widehat H_{ij}(x)|\le C(1+|x|)^{-6}.
\tag{2.4}
\]

This argument avoids assuming a sixth classical derivative at the
origin. The smooth part away from zero decays arbitrarily rapidly.
Twisted finite-torus kernels are exact phase-weighted periodizations of
(2.4). Since the exponent six exceeds the dimension four,

\[
 \boxed{\quad
 |\operatorname{Cov}(\sqrt\beta F_{e;\mu}^{(1)},
                    \sqrt\beta F_{d;\rho}^{(1)}\mid\theta)|
 \le C(1+\operatorname{dist}_{\mathbb T_n}(e,d))^{-6}.
 \quad}
\tag{2.5}
\]

The estimate is uniform through twisted near-zero modes; harmonic modes
contribute zero. Fixed color multiplicities, root offsets and pointing
completions only change C.

The leading reference ordinary root is

\[
 r_e^{(0)}=\beta\sum_{\mu,A}|F_{e,A;\mu}^{(1)}|^2
                  -\ell_e^{(0)}/3.
\]

Its Casimir/Hessian term is deterministic conditional on the flat
background. Wick's formula gives, for uncompleted roots,

\[
 \operatorname{Cov}(r_e^{(0)},r_d^{(0)}\mid\theta)
 =2\sum_{\mu,\rho,A,B}
 \operatorname{Cov}(\sqrt\beta F_{e,A;\mu}^{(1)},
                    \sqrt\beta F_{d,B;\rho}^{(1)}\mid\theta)^2.
\tag{2.6}
\]

In a real mode basis this formula counts conjugate colors and momenta
once. It implies

\[
 |\operatorname{Cov}(r_e^{(0)},r_d^{(0)}\mid\theta)|
 \le C(1+\operatorname{dist}(e,d))^{-12}.
\tag{2.7}
\]

The same bound holds for the fixed finite completion by summing its
literal shifted terms.

\hypertarget{shared-flat-holonomy-means-remain-in-the-bound}{%
\subsubsection{Shared flat-holonomy means remain in the
bound}\label{shared-flat-holonomy-means-remain-in-the-bound}}

The variance identity gives that the per-anchor reference ordinary mean satisfies

\[
 |m_e(\theta)-m_e(0)|\le Cn^{-6}.
\tag{2.8}
\]

For completeness, its symbol is of order \textbar k\textbar\^{}2 at
zero, so the same annular proof gives Fourier coefficients
O((1+\textbar z\textbar)\^{}-6). The exact shifted-grid Poisson identity
then gives a total mean change O(n\^{}-2), or O(n\^{}-6) per anchor. No
common zero mode is discarded.

Under an arbitrary source-independent probability law pi of commuting
flat holonomies,

\[
 |\operatorname{Cov}_{\pi}(m_e(\theta),m_d(\theta))|
 \le Cn^{-12}.
\]

Combining with the conditional covariance and noting that torus
distances are O(n),

\[
 \boxed{\quad
 |\operatorname{Cov}_{\pi(d\theta)\mu_\theta(dA)}
        (r_e^{(0)},r_d^{(0)})|
 \le C(1+\operatorname{dist}(e,d))^{-12}.
 \quad}
\tag{2.9}
\]

This strengthens the finite-volume Schur bound by exhibiting its spatial
kernel. In particular the far covariance tail obeys

\[
 \sup_e\sum_{d:\operatorname{dist}(e,d)\ge R}
 |\operatorname{Cov}(r_e^{(0)},r_d^{(0)})|
 \le C(1+R)^{-8}.
\tag{2.10}
\]

For disjoint marked supports at distance at least R, Schur's test gives

\[
 |\operatorname{Cov}(R_f^{(0)},R_g^{(0)})|
 \le C(1+R)^{-8}\sqrt{M_fM_g}.
\tag{2.11}
\]

This is a statement about the reference source before an arbitrary
observation. A future observation preserves the complete L2 operator
bound but need not preserve (2.9) or (2.10); Section 4 gives an exact
reason.

\hypertarget{source-specific-low-spectral-and-low-mode-concentration-estimates}{%
\subsection{Source-specific low-spectral and low-mode concentration
estimates}\label{source-specific-low-spectral-and-low-mode-concentration-estimates}}

\hypertarget{low-ornstein-uhlenbeck-spectrum}{%
\subsubsection{Low Ornstein-Uhlenbeck
spectrum}\label{low-ornstein-uhlenbeck-spectrum}}

For fixed theta, whiten the nonzero transverse Gaussian modes. Let
L\_theta be the positive Ornstein-Uhlenbeck generator whose one-particle
rates are s(k). The fixed upsilon factor may equivalently be retained in
those rates; it changes only constants. On a centered degree-two mode
indexed by k,q, the rate is s(k)+s(q).

Set

\[
 F_f^{\theta}=R_f^{(0)}-E_\theta R_f^{(0)}.
\]

In whitened real coordinates Z this is

\[
 F_f^{\theta}=Z^TA_fZ-\operatorname{tr}A_f.
\]

The matrix entries in Fourier coordinates satisfy

\[
 \|A_f(k,q)\|_{\rm HS}^2
 \le\frac{C}{V^2}s(k)s(q)
             \sum_\nu|\widehat f_\nu(k-q)|^2.
\tag{3.1}
\]

This follows directly from (2.1): each whitened directional force row
has norm bounded by C sqrt(s(k))/sqrt(V). The finite color and
orientation multiplicity is included in C. In a flat adjoint color
sector the twists cancel in the difference grid. Real modes give the
same estimate with fixed multiplicities.

There are at most C(1+V epsilon\^{}2) shifted-grid momenta with s(k)
\(\leq\) epsilon for 0\textless epsilon \(\leq\) 1. On the low
two-particle spectrum, s(k)s(q) \(\leq\) epsilon\^{}2/4. Grouping by k-q
and using Parseval therefore gives

\[
 \boxed{\quad
 \|1_{(0,\varepsilon]}(L_\theta)F_f^{\theta}\|_2^2
 \le C\bigl(\varepsilon^4+\varepsilon^2n^{-4}\bigr)M_f,
 \qquad 0<\varepsilon\le1.
 \quad}
\tag{3.2}
\]

In detail, the sum over k for each fixed difference is bounded by
\(C\varepsilon^2(1+V\varepsilon^2)\), while
\(\sum_m|\widehat f_\nu(m)|^2=V\sum_x|f_{x,\nu}|^2\). The real Gaussian
variance is twice the relevant Hilbert-Schmidt square.

For the untwisted grid, the source annihilates harmonic modes and the
first positive one-particle rate is bounded below by c n\^{}-2. If the
cutoff reaches positive source spectrum, epsilon \(\geq\) c'n\^{}-2,
which absorbs the second term. Thus

\[
 \boxed{\quad
 \|1_{(0,\varepsilon]}(L_0)F_f^{0}\|_2^2
 \le C\varepsilon^4M_f.
 \quad}
\tag{3.3}
\]

The finite-volume correction in (3.2) is necessary uniformly over
twists: a nonzero shifted momentum can approach zero much faster than
n\^{}-1. Its quadratic source amplitude still vanishes, but the number
of such modes is discrete rather than V epsilon\^{}2.

Under an arbitrary holonomy mixture, let the generator act fiberwise and
leave theta fixed. Its zero eigenspace includes every function of theta.
The globally centered source has the orthogonal decomposition

\[
 R_f^{(0)}-ER_f^{(0)}
 =F_f^{\theta}+[m_f(\theta)-E_\pi m_f].
\]

Equation (2.8) and the original count O(n\^{}4) give

\[
 \operatorname{Var}_\pi m_f(\theta)\le Cn^{-8}M_f.
\tag{3.4}
\]

Therefore the inclusive low projection, with the shared theta mean
retained, obeys

\[
 \boxed{\quad
 \|1_{[0,\varepsilon]}(L)
          (R_f^{(0)}-ER_f^{(0)})\|_2^2
 \le C(\varepsilon^4+n^{-8})M_f.
 \quad}
\tag{3.5}
\]

No zero-energy history mode is removed by a centering convention. The
simple inequality \(\varepsilon^2n^{-4}\le(\varepsilon^4+n^{-8})/2\) was
used in (3.5).

The power four is sharp for the uniform ordinary source in the regime n
sqrt(epsilon) large. The variance identity gives \(K^2/3\preceq T_D\preceq K^2\).
Hence, on transverse modes, the whitened uniform-source quadratic matrix
has eigenvalues comparable to s(k). Its low variance is comparable to

\[
 \sum_{0<s(k)\le\varepsilon/2}s(k)^2
 \asymp V\varepsilon^4
\]

when a fixed small momentum annulus contains many grid points. The
exponent is the four-dimensional phase-space power epsilon\^{}2 times
the squared differentiated-source amplitude epsilon\^{}2.

\hypertarget{a-genuine-reference-low-mode-observation}{%
\subsubsection{A genuine reference low-mode
observation}\label{a-genuine-reference-low-mode-observation}}

Let G\_epsilon be the sigma-algebra generated by theta and the field
modes with s(k) \(\leq\) epsilon. It is a complete positive observation
of this same reference probability law. Define

\[
 G_f=E[R_f^{(0)}|\mathcal G_\varepsilon]-ER_f^{(0)}.
\]

Conditional Gaussian integration leaves the low-low quadratic part, plus
the full mean m\_f(theta). Its centered quadratic part lies in
two-particle spectrum at most 2 epsilon. Consequently

\[
 \|G_f\|_2^2\le C(\varepsilon^4+n^{-8})M_f.
\tag{3.6}
\]

For an observation retaining wavelengths at least R, epsilon is of order
R\^{}-2, giving

\[
 \|G_f\|_2^2\le C(R^{-8}+n^{-8})M_f.
\tag{3.7}
\]

Any later positive observation or orthogonal output projection of this
same reference law contracts this complete norm. There is no
identification of G\_epsilon with the original Wilson X, O, B, or an
actual R-II output.

\hypertarget{concentration-including-the-shared-holonomy-mean}{%
\subsubsection{Concentration, including the shared holonomy
mean}\label{concentration-including-the-shared-holonomy-mean}}

Let b\_f=max\_e\textbar f\_e\textbar. Write the retained centered
quadratic matrix as
\(A_{f,\varepsilon}=P_\varepsilon A_fP_\varepsilon\). For the positive
local root forms B\_e of the finite-volume calculation, \(\sum_e B_e\preceq CK\). Thus

\[
 -b_f\sum_eB_e\preceq A_f\preceq b_f\sum_eB_e,
 \qquad
 \|A_{f,\varepsilon}\|\le Cb_f\varepsilon.
\tag{3.8}
\]

The preceding Fourier count also gives

\[
 \operatorname{tr}(A_{f,\varepsilon}^2)
 \le C(\varepsilon^4+\varepsilon^2n^{-4})M_f.
\tag{3.9}
\]

For a real Gaussian quadratic form the exact determinant formula and its
convergent series imply

\[
 \log E_\theta e^{t(Z^TAZ-\operatorname{tr}A)}
 \le\frac{t^2\operatorname{tr}(A^2)}{1-2|t|\|A\|}.
\tag{3.10}
\]

Meanwhile (2.8) bounds the range of m\_f(theta) by

\[
 Cn^{-6}\sum_e|f_e|
 \le Cn^{-4}\sqrt{M_f}.
\]

Hoeffding's bounded-variable estimate, applied only to this already
identified shared mean, supplies a subgaussian log-mgf C n\^{}-8 M\_f
t\^{}2. The conditional quadratic estimate is uniform in theta.
Combining the two gives

\[
 \boxed{\quad
 \log E e^{tG_f}
 \le\frac{C M_f(\varepsilon^4+n^{-8})t^2}
               {2(1-C b_f\varepsilon|t|)},
 \qquad Cb_f\varepsilon|t|<1.
 \quad}
\tag{3.11}
\]

This estimate is for the actual reference mixture, with arbitrary
source-independent pi and the complete shared mean retained. It has no
beta-dependent loss. For a random weight 0 \(\leq\) H \(\leq\) 1 of
probability p, the standard tail integration now yields

\[
 \boxed{\quad
 E[HG_f^2]
 \le Cp\left[
 M_f(\varepsilon^4+n^{-8})\log(e/p)
 +b_f^2\varepsilon^2\log^2(e/p)
 \right].
 \quad}
\tag{3.12}
\]

Conditional Jensen preserves (3.11) under any subsequent positive
observation of the same reference law. Orthogonal output projections
preserve the corresponding L2 bound on the complete variable or already
formed rare lineage, not an arbitrary newly masked tail.

The exact physical squared factor \(c_\beta^{-2}a^{-4}\) remains outside
all dimensionless variance estimates. No old source grade, coupling
ratio, age factor, raw parent mass, or random history coefficient has
been absorbed into C. Such coefficients must be reinserted with their
actual deterministic status or separately proved envelope.

\hypertarget{what-a-nonlinear-upgrade-would-have-to-prove}{%
\subsubsection{What a nonlinear upgrade would have to
prove}\label{what-a-nonlinear-upgrade-would-have-to-prove}}

The reference calculation exhibits a source-specific infrared mechanism:
two differentiated curvature legs and four-dimensional phase space
suppress low energy by epsilon\^{}4. It does not use a global mass gap.

For the actual ordinary retained source and its actual marginal
generator B\_X, a sufficient analogue would be a uniform source-specific
estimate such as

\[
 \|1_{(0,\varepsilon]}(B_X)R_{X,f,c}\|_2^2
 \le C\varepsilon^\alpha M_f
\]

with an explicitly controlled shared kernel component and
alpha\textgreater0, together with the existing high-spectrum bound.
Equivalently, an appropriate annealed semigroup correlation estimate
would directly supply the missing low-spectral control. No such
nonlinear bound follows from (3.2), from the local moments of the odd-torus transfer theorem, or from
one-link conditional gaps. The induced density after a color integration
includes its raw normalizers and effective scores, and the actual sharp
domains need not be convex. This note proves no comparison of that
density or generator with L\_theta.

\hypertarget{an-actual-observation-can-make-conditional-mean-covariance-long-range}{%
\subsection{An actual observation can make conditional-mean covariance
long
range}\label{an-actual-observation-can-make-conditional-mean-covariance-long-range}}

This is an exact probability identity applicable to the original Wilson
law, not only the Gaussian benchmark.

Let r\_x be the complete ordinary source summed over the finitely many
orientations/pointed members in one original lattice cell. The original
Wilson law and the completed source are translation covariant. Put

\[
 F=\sum_{x\in\mathbb T_n^4}r_x,
 \qquad Y=F.
\]

The deterministic kernel reading Y is a positive source-independent
observation: it reads the original source-zero field and changes no
measure. By translation invariance, E(r\_x\textbar F) is independent of
x. Summing over x gives exactly

\[
 E(r_x|F)=F/V,
 \qquad
 \operatorname{Cov}(E[r_x|F],E[r_y|F])
 =\operatorname{Var}(F)/V^2
\tag{4.1}
\]

for all x,y, including maximally separated ones. Whenever the
original-law extensive lower bound applies,
\(\operatorname{Var}(F)\ge cV\), this conditional-mean covariance is at
least c/V everywhere. In the Gaussian benchmark the matching upper bound
also holds, so it is of order V\^{}-1.

At distances of order n, V\textsuperscript{-1=n}-4 is much larger than
n\^{}-12. Therefore arbitrary positive observation does not preserve the
spatial covariance decay in (2.9), even though its complete covariance
operator remains bounded whenever the original one is bounded. The
resulting rank-one covariance matrix has eigenvalue Var(F)/V, which is
of order one under an extensive susceptibility bound.

This observation is not claimed to be an R-II history and is not a
counterexample to (11). It proves that spatial localization of
conditional means requires actual restrictions on the history/readout;
Jensen alone gives only the complete marked L2 bound. It also explains
why future rootwise masks or deletion of a shared history mean can
change the question.

\hypertarget{conclusion-of-part-iii}{%
\subsection{Conclusion of
Part III}\label{conclusion-of-part-iii}}

The original-law O/B construction, if its full-factor hypotheses are
verified, supplies a bounded local covariance subtraction and an exact
equivalence of susceptibility bounds. Its retained covariance remains
the unknown quantity.

The flat benchmark now has a quantitative source-specific infrared
exponent, a spatial covariance kernel, and a low-mode concentration
theorem retaining flat-holonomy history means. These establish the
mechanism one would want a nonlinear renormalization argument to
preserve. They do not establish that preservation, mixed-domain
admission, or S94.

The actual Wilson ordinary estimate (11) is therefore still open. No
local moment or conditional Gaussian residual has been used as a
replacement for the missing spatial/low-spectral theorem.

\newpage

\hypertarget{part-iv.-infrared-bounds-for-the-full-compact-quadratic-jet-in-a-reference-law}{%
\section{Part IV. Infrared bounds for the full compact quadratic jet in
a reference
law}\label{part-iv.-infrared-bounds-for-the-full-compact-quadratic-jet-in-a-reference-law}}

\hypertarget{scope}{%
\subsection{Scope}\label{scope}}

This note combines the actual fixed-SU(N) compact conditional-Hessian
result proved by Part II with an exact Gaussian calculation. The rank N
is fixed and all group factors remain in the constants. It does not
replace the compact marginal law by a Gaussian. The observable used
below is the complete field-degree-two jet of the actual compact
conditional ordinary source at the flat field, evaluated under an
explicitly named Gaussian marginal reference.

The calculation keeps the O(1) lower-derivative Hessian remainder,
including the full Casimir and compact-normalizer effects. After
canonical normalization this is O(beta\^{}-1); it permits a lower
infrared power than the pure leading-source epsilon\^{}4 law. The
resulting mixed powers are still summable. No arbitrary flat-holonomy
mixture estimate for this additional remainder is claimed.

\hypertarget{actual-compact-input-with-local-source-structure}{%
\subsection{Actual compact input, with local source
structure}\label{actual-compact-input-with-local-source-structure}}

Let X retain all links except a fixed periodic separated set of seed
links. The period cell and all original source determining collars are
fixed independently of n and beta. Let H0 be the full dimensionless flat
Wilson/Maxwell Hessian, with blocks indexed by eliminated x and retained
y:

\[
 H_0=\begin{pmatrix}H_{xx}&H_{xy}\\H_{yx}&H_{yy}\end{pmatrix},
 \quad
 E=\binom{-H_{xx}^{-1}H_{xy}}{I},
 \quad
 S=H_{yy}-H_{yx}H_{xx}^{-1}H_{xy}.
\tag{1.1}
\]

H0 is the physical flat Hessian on the gauge/zero-mode quotient. The
actual compact induced action is defined relative to the ORIGINAL
retained product Haar measure. No retained-coordinate Haar Jacobian is
absorbed into an arbitrary Lebesgue-density Hessian; such a coordinate
Hessian can add O(1) terms and gauge artifacts. H\_xx is a fixed
positive diagonal block because no original plaquette reads two
eliminated seeds. In the unit Maxwell convention its link coefficient is
six. The exact identities are

\[
 H_0E=\binom0S,
 \qquad E^*H_0E=S,
 \qquad E^*H_0^2E=S^2.
\tag{1.2}
\]

All maps in (1.1) are finite range and uniformly bounded, because the
eliminated block is diagonal with a fixed lower bound.

For a complete original ordinary root r\_c, set
m\_c=E\_mu(r\_c\textbar X). The induced-layer result proves at the flat
field

\[
 \operatorname{Hess}m_c
 =\beta E^*K_{D,c}E+E_{\beta,c},
\tag{1.3}
\]

where the leading form K\_D,c is the original complete
directional-square Hessian, with fixed normalization factors included.
The error E\_beta,c has fixed support and uniformly O(1) coefficients
for beta \(\geq\) 1. It is gauge invariant, has zero first field jet at
flat, and annihilates every local gradient. Its source-normalizer proof
retains the actual compact denominator.

A fixed contractible collar therefore provides a cycle factorization

\[
 E_{\beta,c}=C_c^*A_{\beta,c}C_c,
 \qquad \|A_{\beta,c}\|\le C.
\tag{1.4}
\]

Here C\_c reads a finite set of retained-link cycles. These cycles can
be filled by ORIGINAL plaquettes in the full collar, including
plaquettes containing eliminated links; one must not require fillings
made only of retained elementary plaquettes. A cycle supported on
retained links has the same value on y and on its harmonic extension Ey.

The finite-dimensional factorization follows because a symmetric local
Hessian annihilating gradients factors through the cycle space. A fixed
cycle basis and fixed plaquette fillings have bounded norms. This also
explains why local flat torons are annihilated: inside a fixed box their
linearized constant field is a gradient, and every contractible cycle
has zero first-order holonomy on it.

Write the leading forms as sums of a fixed number of local
directional-force rows L\_c:

\[
 E^*K_{D,c}E=L_c^*L_c.
\]

The complete flat inequality K\_D \(\leq\) C H0\^{}2 and the fixed cycle
fillings give

\[
 \sum_cL_c^*L_c\preceq CS^2,
 \qquad
 \sum_cC_c^*C_c\preceq CS.
\tag{1.5}
\]

The second inequality is bounded-overlap local filling followed by
E\emph{H0E=S. The first uses the exact E}H0\textsuperscript{2E=S}2; no
smallest positive eigenvalue is divided out.

For deterministic marks f\_c, the actual compact conditional Hessian is
consequently

\[
 H_f=\sum_cf_c\bigl(\beta L_c^*L_c+C_c^*A_{\beta,c}C_c\bigr).
\tag{1.6}
\]

All original source offsets and tests are included in the finite row
lists. The terms linear in the source amplitude are being retained;
higher external-source grades are separate obligations.

\hypertarget{the-explicitly-specified-gaussian-reference}{%
\subsection{The explicitly specified Gaussian
reference}\label{the-explicitly-specified-gaussian-reference}}

Put a centered Gaussian law on the positive subspace of S with
covariance

\[
 \operatorname{Cov}(y)=(\beta S)^+.
\tag{2.1}
\]

Gauge and exact harmonic zero modes may be fixed or given any law, since
(1.5) shows that the quadratic source annihilates them. Define the
centered quadratic jet

\[
 Q_f(y)=\frac12 y^*H_fy-E\left[\frac12 y^*H_fy\right].
\tag{2.2}
\]

This is a Gaussian benchmark of the actual compact Hessian. It is not
the random variable m\_f(y) under its actual compact marginal law.

Whiten y=beta\^{}-1/2 S\^{}+1/2 Z. The centered quadratic matrix is

\[
 A_f=\frac12\sum_cf_c\left(
 B_c^*B_c+\beta^{-1}D_c^*A_{\beta,c}D_c\right),
 \quad B_c=L_cS^{+1/2},\quad D_c=C_cS^{+1/2}.
\tag{2.3}
\]

Equation (1.5) gives

\[
 \sum_cB_c^*B_c\preceq CS,
 \qquad \sum_cD_c^*D_c\preceq CI
\tag{2.4}
\]

on the positive subspace. The lower-derivative term is explicit and
survives as beta\^{}-1 in (2.3).

\hypertarget{uniform-low-mode-count-for-the-fixed-cell-schur-complement}{%
\subsection{Uniform low-mode count for the fixed-cell Schur
complement}\label{uniform-low-mode-count-for-the-fixed-cell-schur-complement}}

The block Gaussian elimination congruence diagonalizes H0 into
diag(Hxx,S). The congruence matrix and its inverse have uniformly
bounded norms because Hxx\^{}-1 and Hxy do. Ordered eigenvalues,
including the common nullity, are therefore comparable up to fixed
constants.

The positive flat Maxwell spectrum in four dimensions has at most C(1+V
epsilon\^{}2) modes below epsilon, where V=n\^{}4. Since Hxx has a fixed
positive lower bound, for a fixed small epsilon the same counting
estimate holds for S:

\[
 N_S(\varepsilon)
 :=\#\{j:0<\lambda_j(S)\le\varepsilon\}
 \le C(1+V\varepsilon^2).
\tag{3.1}
\]

The first positive eigenvalue of S is bounded below by c n\^{}-2 in the
untwisted periodic model. Gauge nullity has not been counted as small
positive spectrum.

Because the seed pattern is periodic with a fixed cell, eigenvectors may
be chosen in Bloch form. Let V\_cell be the number of period cells, so
V\_cell is a fixed multiple of V. There are only a fixed number of bands
and source types per cell. For a normalized Bloch eigenmode psi\_j(k) of
rate lambda\_j(k), translation covariance and (2.4) imply the per-cell
row estimates

\[
 \|B_c\psi_j(k)\|\le C\sqrt{\lambda_j(k)}/\sqrt{V_{\rm cell}},
 \qquad
 \|D_c\psi_j(k)\|\le C/\sqrt{V_{\rm cell}}.
\tag{3.2}
\]

These bounds do not require smooth choices of eigenvectors through
degenerate gauge bands. They follow from summing the squared row
amplitude over all translated cells.

\hypertarget{the-mixed-power-low-spectral-variance-estimate}{%
\subsection{The mixed-power low-spectral variance
estimate}\label{the-mixed-power-low-spectral-variance-estimate}}

Let L\_S be the positive Gaussian OU generator whose one-particle rates
are the positive eigenvalues of S. A centered quadratic mode has rate
lambda\_j(k)+lambda\_l(q).

For the marked quadratic matrix, (3.2) gives

\[
 \|A_f(k,j;q,l)\|^2
 \le\frac C{V_{\rm cell}^2}
 \sum_\alpha|\widehat f_\alpha(k-q)|^2
 \left(\sqrt{\lambda_j(k)\lambda_l(q)}+\beta^{-1}\right)^2.
\tag{4.1}
\]

The finite index alpha records the original source types within one
period cell, including all pointing offsets. Group by k-q. The low
two-particle spectrum has both rates at most epsilon; for each
difference its number of contributing pairs is at most C(1+V
epsilon\^{}2), by (3.1) and the fixed band count. Parseval and the real
Gaussian quadratic-variance formula then yield

\[
 \|1_{(0,\varepsilon]}(L_S)Q_f\|_2^2
 \le C M_f(\varepsilon+\beta^{-1})^2
                 (\varepsilon^2+V^{-1}).
\tag{4.2}
\]

In the untwisted model, positive source spectrum is absent below c
n\^{}-2. Once the cutoff reaches it, V\^{}-1 \(\leq\) C epsilon\^{}2.
Thus

\[
 \boxed{\quad
 \|1_{(0,\varepsilon]}(L_S)Q_f\|_2^2
 \le C M_f\varepsilon^2(\varepsilon+\beta^{-1})^2
 \le C M_f\left(
 \varepsilon^4+\beta^{-1}\varepsilon^3
                  +\beta^{-2}\varepsilon^2\right),
 \quad 0<\varepsilon\le1.
 \quad}
\tag{4.3}
\]

Constants in the last expression absorb the exact cross coefficient.
Equivalently one may use C M\_f(epsilon\textsuperscript{4+beta}-2
epsilon\^{}2) by the elementary square inequality. Keeping all three
powers makes the source mixing visible.

If a Gaussian observation retains precisely the S-modes of rate at most
epsilon, conditional expectation of Q\_f retains its low-low quadratic
part. The same estimate holds, up to replacing epsilon by 2 epsilon.
Later positive observations and orthogonal projections of this same
Gaussian law contract that complete bound.

In length-scale notation epsilon=R\^{}-2, the result is

\[
 \operatorname{Var}(E[Q_f|\text{modes }\lambda\le R^{-2}])
 \le C M_f\left(R^{-8}+\beta^{-1}R^{-6}
                           +\beta^{-2}R^{-4}\right).
\tag{4.4}
\]

This is the correct benchmark for the actual compact conditional jet.
The pure leading-source R\^{}-8 law is not asserted for its full jet. A
nonzero lower-derivative remainder can dominate when R\^{}2 is much
larger than beta. The bound permits that term; it does not prove that
its coefficient is nonzero for every particular source.

\hypertarget{matching-gaussian-concentration}{%
\subsection{Matching Gaussian
concentration}\label{matching-gaussian-concentration}}

Let P\_epsilon be the one-particle low-mode projection. From (2.4), for
b\_f=max\textbar f\_c\textbar,

\[
 \|P_\varepsilon A_fP_\varepsilon\|
 \le Cb_f(\varepsilon+\beta^{-1}).
\tag{5.1}
\]

The Hilbert-Schmidt square is bounded by the right side of (4.3). The
exact centered Gaussian quadratic determinant formula therefore gives,
for the low-mode conditional quadratic G\_f,

\[
 \boxed{\quad
 \log E e^{tG_f}
 \le\frac{C M_f\varepsilon^2(\varepsilon+\beta^{-1})^2 t^2}
              {2[1-Cb_f(\varepsilon+\beta^{-1})|t|]},
 \qquad Cb_f(\varepsilon+\beta^{-1})|t|<1.
 \quad}
\tag{5.2}
\]

Here G\_f is centered by its complete mean in the named Gaussian
reference. No actual random compact-history mean has been removed. For a
rare weight H with EH=p, this also gives

\[
 E[HG_f^2]\le Cp\left[
 M_f\varepsilon^2(\varepsilon+\beta^{-1})^2\log(e/p)
 +b_f^2(\varepsilon+\beta^{-1})^2\log^2(e/p)
 \right].
\tag{5.3}
\]

This remains a reference result. It supplies no moment or zero-free
source estimate for the nonlinear compact conditional density.

\hypertarget{spatial-covariance-of-the-same-compact-jet-benchmark}{%
\subsection{Spatial covariance of the same compact-jet
benchmark}\label{spatial-covariance-of-the-same-compact-jet-benchmark}}

There is a matching spatial estimate. Let the root quadratic jet be
Q\_c.~The local leading force rows L\_c annihilate affine link fields:
the harmonic extension E reproduces an affine field because H0
annihilates it locally, and each directional second-difference row
annihilates it as well. Their coefficients thus have vanishing moments
of orders zero and one. The cycle rows C\_c annihilate constant link
fields, so their coefficients have zero zeroth moment.

For any retained gauge-invariant linear rows w,v, let their full-link
extensions be tilde w=(0,w), tilde v=(0,v). They are orthogonal to ker
H0, and the Schur identity gives

\[
 wS^+v^*=\widetilde w H_0^+\widetilde v^*.
\tag{6.1}
\]

One proof solves H0 z=tilde v: its retained component solves S z\_y=v,
and pairing with w removes the arbitrary gauge solution. Since the
extended rows annihilate gradients, the transverse projector in H0\^{}+
drops out. Their covariance is a local differentiated scalar lattice
Green function.

A dyadic Fourier estimate on the finite torus gives uniform decay powers
determined by the vanishing moments. For two cycle rows the multiplier
has order zero at k=0 and its kernel decays as distance\^{}-4. A
leading/cycle pair has order one and distance\^{}-5. Two leading rows
have order two and distance\^{}-6. Exact zero modes contribute zero.
These finite-torus estimates can be proved directly by summing dyadic
annuli down to scale n\^{}-1; the lowest annulus contributes at most
Cn\^{}-4, Cn\^{}-5, or Cn\^{}-6, respectively, and no absolute
periodization at the borderline exponent four is assumed.

Consequently, in the reference (2.1), Wick contraction gives

\[
 \boxed{\quad
 |\operatorname{Cov}(Q_c,Q_d)|
 \le C\left[
 (1+\operatorname{dist}(c,d))^{-12}
 +\beta^{-1}(1+\operatorname{dist}(c,d))^{-10}
 +\beta^{-2}(1+\operatorname{dist}(c,d))^{-8}
 \right].
 \quad}
\tag{6.2}
\]

The signed local matrices A\_beta,c in the lower-order term cause no
difficulty: their norms and ranks are uniformly bounded, and the
absolute Wick cross terms use the corresponding linear-row covariance
bounds. Fixed cell types and pointing offsets change only constants.

The far covariance row sum is therefore

\[
 \sup_c\sum_{d:\operatorname{dist}(c,d)\ge R}
 |\operatorname{Cov}(Q_c,Q_d)|
 \le C\left(R^{-8}+\beta^{-1}R^{-6}
                         +\beta^{-2}R^{-4}\right)
\tag{6.3}
\]

for R \(\geq\) 1. This agrees with the low-spectral power count. All
three pieces are summable in four spatial dimensions.

\hypertarget{exact-limitation-of-the-conclusion}{%
\subsection{Exact limitation of the
conclusion}\label{exact-limitation-of-the-conclusion}}

The compact calculation proves the actual source's flat Hessian
decomposition (1.3), including its O(1) local remainder. The Gaussian
calculation proves (4.3), (5.2), and (6.2) for that Hessian as a
quadratic observable under the explicitly specified reference law.

To obtain the original Wilson retained-ordinary estimate, one still
needs a uniform nonlinear comparison controlling the entire actual
source, its source-dependent induced density, the random background and
history means, and the full projected cross terms. The actual sharp
domains and mixed-history normalizers are not replaced by (2.1). The
reference generator L\_S is not the original-law compressed heat-bath
generator A\_B.

The physical squared factor c\_beta\^{}-2 a\^{}-4, original source
grades, age and coupling-ratio factors, and parent/selected scalar
masses remain external to this dimensionless benchmark. No source
redefinition or discarded Casimir constant is involved.

\newpage

\hypertarget{part-v.-actual-nonlinear-remainder-and-the-remaining-spatial-count}{%
\section{Part V. Actual nonlinear remainder and the remaining spatial
count}\label{part-v.-actual-nonlinear-remainder-and-the-remaining-spatial-count}}

\hypertarget{scope-and-theorem}{%
\subsection{Scope and theorem}\label{scope-and-theorem}}

This result concerns the fixed complete compact observation X of the
original Wilson law, obtained by integrating a fixed separated set of
original links. It does not concern a sharply selected good/bad fiber,
and it does not identify X with B*, R-II, or a later mixed-history
readout.

Keep the original complete ordinary source

\[
 R_f=\sum_cf_c r_c^0,\qquad
 r_c^0=\beta d_c-\ell_c/3,
 \qquad M_f=\sum_c|f_c|^2,
\]

with the full color sums, directional squares, neighboring roots,
pointed offsets and original Casimir response. Marks f\_c are
deterministic original marks. Define

\[
 m_c=P_Xr_c^0.
\]

The original multipoint/local Wilson moment input is used only on fixed
literal plaquette records. The rank N, source stencils, separated-link
pattern, and all local collars remain fixed. Constants do not depend on
beta \(\geq\) 1 or the torus size once the collars embed and the stated
moment regime applies.

For each root c one can construct, using only observed links, a fixed
local tree gauge and a smooth local cutoff chi\_c.~Let y\_c denote the
resulting real retained-link logarithm coordinates on that chart, and
let

\[
 H_c=D_y^2m_c(0),\qquad
 q_c=\frac12\chi_c H_c[y_c,y_c],\qquad
 e_c=m_c-m_c(0)-q_c.
\tag{1.1}
\]

Outside the chart q\_c is defined to be zero. These are all X-measurable
local functions. H\_c is the actual complete compact conditional
Hessian: all derivatives of the genuine denominator and all readout
responses remain. In particular its O(1) lower-derivative remainder and
full Casimir contribution are included.

The theorem is

\[
 \boxed{\quad
 \|e_c\|_{L^q(\mu_\beta)}
 \le C\beta^{-1/2}q^{3/2},\qquad q\ge2.
 \quad}
\tag{1.2}
\]

For the centered complete error

\[
 E_f=\sum_cf_c(e_c-Ee_c),
\]

the actual original-link heat-bath Dirichlet form obeys

\[
 \boxed{\quad
 \mathcal E_{\rm orig}(E_f)\le C\beta^{-1}M_f.
 \quad}
\tag{1.3}
\]

Let A\_X=P\_X A\_orig P\_X restricted to range(P\_X), with the actual
original-law marginal and all induced normalizers. Then

\[
 \boxed{\quad
 \|1_{[\delta,\infty)}(A_X)E_f\|_2^2
 \le\frac{C}{\beta\delta}M_f,\qquad\delta>0.
 \quad}
\tag{1.4}
\]

The unconditional variance estimate supplied by these inputs is instead

\[
 \boxed{\quad
 \|E_f\|_2^2\le C\beta^{-1}
             \left(\sum_c|f_c|\right)^2.
 \quad}
\tag{1.5}
\]

Equation (1.5) is deliberately not replaced by C beta\^{}-1 M\_f.~The
difference is the exact unresolved spatial count. For a uniform mark on
V roots it is C V\^{}2/beta. It is extensive by this estimate only in
the restricted regime V \(\lesssim\) beta. The quadratic q\_c is
evaluated in actual original-law chart variables; it is not a Gaussian
random variable.

\hypertarget{the-exact-compact-conditional-means-and-their-denominators}{%
\subsection{The exact compact conditional means and their
denominators}\label{the-exact-compact-conditional-means-and-their-denominators}}

Let S0 be the fixed set of eliminated original links. Separation ensures
that no original action plaquette reads two links in S0. Conditional on
all observed links, their compact Wilson kernels factor. A complete
source root reads only a fixed finite number of these kernels, because
its original determining stencil is finite.

For the affected free links, write the normalized conditional integral
schematically as

\[
 m_c(y)=
 \frac{\int r_c^0(u,y)e^{-\beta S_{\rm loc}(u,y)}\,d\lambda(u)}
      {\int e^{-\beta S_{\rm loc}(u,y)}\,d\lambda(u)}.
\tag{2.1}
\]

All actual exterior scalar factors and every compact normalizer remain
in the raw X marginal. Normalizers of independent free links not read by
r\_c cancel in this normalized mean. The denominator of each affected
kernel stays in (2.1). Derivatives of m\_c therefore differentiate its
real numerator, denominator, material readout and every affected
neighboring/pointed source member.

The fixed local gauge below is used to evaluate a gauge-invariant
function and to estimate its moments. It is not a change to a
gauge-fixed probability density. No Faddeev-Popov factor, new reference
law, or dropped Haar Jacobian is introduced.

\hypertarget{an-observed-link-spanning-tree-and-a-fixed-plaquette-record}{%
\subsection{An observed-link spanning tree and a fixed plaquette
record}\label{an-observed-link-spanning-tree-and-a-fixed-plaquette-record}}

Take a fixed original box containing the complete determining stencil of
r\_c and the full action stars of the finitely many eliminated links it
reads. Enlarge it by a fixed number of lattice units as needed.

Start with a fixed spanning tree in the original box. Whenever a tree
edge is eliminated, replace it by the other three edges of an elementary
plaquette through that edge. Those three edges are observed, because no
action plaquette contains two eliminated links. The union of these
detours is a connected observed-edge graph. Choose a deterministic
spanning tree from it and the needed observed determining edges. All
choices depend only on the fixed lattice/seed geometry, not on U. The
enlarged collar and the number of edges remain bounded independently of
n and beta.

Products of observed links along this tree define a local gauge
transformation measurable with respect to X. Every tree edge becomes the
identity. Each remaining observed edge in the determining collar becomes
a rooted loop holonomy made of that edge and tree paths. Each such loop
lies in the fixed box and admits a fixed filling by original plaquettes.
Fillings may contain eliminated links; the LOOP ITSELF and the tree use
observed links only.

Let P\_c be the finite union of those original plaquette fillings, with
all repeated incidences counted by a fixed multiplicity. A nonabelian
filling identity expresses each loop holonomy as a finite product of
conjugated original plaquette holonomies and inverses. Bi-invariant
distance and the triangle inequality give

\[
 \max_j d(I,V_{c,j})^2
 \le C\sum_{p\in P_c}s_p,
\tag{3.1}
\]

where V\_c,j are the observed tree-gauged links and
\(s_p=N-\operatorname{Re}\operatorname{Tr}U_p\). One may also use the
exact identity \(\|I-U_p\|_{\rm HS}^2=2s_p\) and telescoping products.
The fixed filling bound is essential; an arbitrary long tree does not
satisfy (3.1).

On a sufficiently small fixed identity chart, set V\_c,j=exp(y\_c,j).
Then

\[
 |y_c|^2\le C\sum_{p\in P_c}s_p.
\tag{3.2}
\]

The coordinates y\_c and chart membership are X-measurable, even though
the original plaquette record used to bound them can contain eliminated
links. Inequality (3.1) holds for every original field realizing that
observed datum, so taking original-law moments is legitimate.

Gauge covariance of the completed root and Haar invariance of the
compact free-link integrals show that m\_c evaluated on the original
bath equals m\_c evaluated on this locally tree-gauged bath. The
residual global conjugation at the tree root does not matter: m\_c, its
Hessian, and the norm cutoffs are invariant under it.

\hypertarget{actual-chart-moments-and-rare-chart-failure}{%
\subsection{Actual chart moments and rare chart
failure}\label{actual-chart-moments-and-rare-chart-failure}}

For the fixed literal record P\_c, the odd-torus transfer theorem supplies constants
theta,K\textgreater0 such that

\[
 E_{\mu_\beta}\exp\left(\theta\beta\sum_{p\in P_c}s_p\right)
 \le K.
\tag{4.1}
\]

The exact admissible n,beta regime of this imported input is retained.
No mixing statement is added to it. Equations (3.1)-(3.2) imply, on the
outer logarithm chart,

\[
 \|\sqrt\beta\,y_c\|_{L^p}\le C\sqrt p,
 \qquad p\ge2,
\tag{4.2}
\]

with the chart indicator understood. They also imply exponentially small
failure of any fixed smaller chart:

\[
 \Pr\{\max_jd(I,V_{c,j})>\delta\}
 \le K e^{-\gamma_\delta\beta},
\tag{4.3}
\]

where gamma\_delta\textgreater0 is fixed. This event is an actual
original-law event; it is not a normalized bad branch.

Choose inner and outer chart radii
0\textless delta\_in\textless delta\_out within the real regularity
chart below. Let chi\_c be one on the inner chart and supported in the
outer chart, with 0 \(\leq\) chi\_c \(\leq\) 1. It can be chosen
smoothly as a function of the tree-gauged logarithm norms. Set q\_c=0
outside its support. The event on which chi\_c differs from one has
probability at most K exp(-gamma beta).

\hypertarget{uniform-real-c3-control-of-the-actual-complete-mean}{%
\subsection{Uniform real C3 control of the actual complete
mean}\label{uniform-real-c3-control-of-the-actual-complete-mean}}

This is the local lemma proved in the induced-layer calculation, Section 9,
equations LC1-LC5. Its proof is included here to make the dependency
explicit.

At the identity bath, each full compact one-link action is 6s(u), up to
the unchanged exterior scalar, and has a unique minimum at u=I with
positive Hessian. Compactness gives a fixed action separation outside an
identity neighborhood. For a sufficiently small fixed REAL bath chart,
the actual action has a unique analytic minimizer u\_*(y), a uniformly
positive local Hessian, and the same fixed separation away from its
minimum.

Use the exact Haar translation u=u\_*(y)v. With the minimum separated
only in this normalized-integral proof, write

\[
 L(v,y)=S_e(u_*(y)v,y)-S_e(u_*(y),y).
\]

The raw factor exp(-beta S\_min(y)) remains in the original marginal
normalizer. It cancels in the ratio defining m\_c.~In a fixed logarithm
chart v=exp(xi),

\[
 L(0,y)=D_\xi L(0,y)=0,
 \qquad L(\xi,y)\ge c|\xi|^2.
\]

Since the first two identities hold for every bath y in the chart, for j
\(\leq\) 3,

\[
 |D_y^j L(\xi,y)|\le C_j|\xi|^2.
\tag{5.1}
\]

The actual normalized compact kernel has Gaussian-scale xi moments.
Hence beta D\_y\^{}j L has uniformly bounded fixed moments. Outside the
free-link chart, every fixed beta power is dominated by the
exponentially small mass, while the true factored compact normalizer is
bounded below by c beta\^{}(-dim(G)/2). For finitely many free links the
same reasoning applies to their fixed-dimensional product.

The complete transported carrier \(A_c(v,y)=r_c^0(u_*(y)v,y)\) and its
first three y derivatives are bounded by C beta on the compact free
fiber. This includes derivatives of u\_*(y), every source readout, and
the original Casimir response. Exact differentiation uses

\[
 D_aK_y(A)=K_y(D_aA)
            -\operatorname{Cov}_{K_y}(A,\beta D_aL).
\tag{5.2}
\]

Repeated differentiation through order three retains every normalizer
term and gives only fixed products of the bounded-moment score
derivatives in (5.1), multiplied by carrier derivatives of size C beta.
Thus

\[
 \sup_{|y|<\delta_{\rm out}}\|D_y^3m_c(y)\|\le C\beta,
 \qquad \|H_c\|\le C\beta.
\tag{5.3}
\]

At flat, D\_y m\_c(0)=0. Gauge/global-adjoint invariance supplies this:
there is no nonzero adjoint-invariant linear functional on the original
link Lie algebra. Also \textbar m\_c(0)\textbar{} \(\leq\) C. For the
latter, on the flat compact kernel the directional forces vanish at the
free-link minimum and grow at most linearly in the local free
coordinate; beta times their squared norm has bounded mean. The Casimir
term is uniformly bounded. Free-link tails are exponentially small. An
atom reading no free link has the same statements directly.

Taylor's theorem therefore gives

\[
 |m_c(y)-m_c(0)-\tfrac12H_c[y,y]|
 \le C\beta|y|^3
\tag{5.4}
\]

on the outer real chart. No complex-bath or sharp selected-domain
regularity is asserted here.

\hypertarget{proof-of-the-actual-remainder-moment-bound}{%
\subsection{Proof of the actual remainder moment
bound}\label{proof-of-the-actual-remainder-moment-bound}}

On the inner chart chi\_c=1, so (5.4) and (4.2) give

\[
 \|1_{\rm inner}e_c\|_q
 \le C\beta\|y_c\|_{3q}^3
 \le C\beta^{-1/2}q^{3/2}.
\tag{6.1}
\]

On the chart-failure/transition event A\_c=\{chi\_c \(\neq\) 1\}, use
the exact complete source moment and conditional Jensen:

\[
 \|m_c\|_{2q}\le\|r_c^0\|_{2q}\le Cq.
\tag{6.2}
\]

The root moment follows from the original local Wilson input, including
its full Casimir constant. Hölder and (4.3) imply

\[
 \|1_{A_c}m_c\|_q\le Cq e^{-\gamma\beta/(2q)}.
\]

The cutoff quadratic is bounded pointwise by C beta on its fixed outer
chart, so

\[
 \|1_{A_c}q_c\|_q\le C\beta e^{-\gamma\beta/q}.
\]

The constant m\_c(0) is uniformly bounded. With x=beta/q, the elementary
bounds on x\textsuperscript{(1/2)e}(-gamma x/2) and
x\textsuperscript{(3/2)e}(-gamma x) show that all these terms are
bounded by

\[
 C\beta^{-1/2}q^{3/2}.
\]

Together with (6.1), this proves (1.2) uniformly for q \(\geq\) 2. No
conditional good-bath probability has been divided out.

The same proof gives \(\|q_c\|_q\le Cq\). The coordinates and cutoff may
differ from root to root, but each uses a fixed original collar and the
source itself has not been changed: equation (1.1) is an exact
decomposition of its actual conditional mean.

\hypertarget{actual-dirichlet-and-high-spectral-consequences}{%
\subsection{Actual Dirichlet and high-spectral
consequences}\label{actual-dirichlet-and-high-spectral-consequences}}

Each e\_c is X-measurable and reads a fixed observed-link collar
L\_c.~The number of such collars through an original link is uniformly
bounded. Let K\_b be the exact original one-link heat bath, and let U'
denote its stationary refresh at b. If b is outside L\_c then Delta\_b
e\_c=0; otherwise stationarity and (1.2) give

\[
 \|e_c(U')-e_c(U)\|_2\le2\|e_c\|_2
 \le C\beta^{-1/2}.
\]

Therefore

\[
 \|\Delta_bE_f\|_2
 \le C\beta^{-1/2}\sum_{c:b\in L_c}|f_c|.
\]

Bounded incidence and Cauchy-Schwarz give

\[
 \mathcal E_{\rm orig}(E_f)
 =\frac12\sum_bE|\Delta_bE_f|^2
 \le C\beta^{-1}M_f.
\]

Because E\_f belongs to range(P\_X),

\[
 \langle E_f,A_XE_f\rangle=\mathcal E_{\rm orig}(E_f).
\]

Spectral calculus now proves (1.4). In particular, the nonlinear
remainder above compressed rate beta\^{}-1 has an extensive bound:

\[
 \|1_{[\beta^{-1},\infty)}(A_X)E_f\|_2^2\le CM_f.
\tag{7.1}
\]

This is a genuine actual-law spectral restriction on the HIGHER-ORDER
ERROR. It does not control the error below beta\^{}-1, nor the quadratic
q\_c under the actual marginal.

\hypertarget{the-genuine-induced-marginal-heat-bath-has-the-same-bound}{%
\subsubsection{The genuine induced marginal heat bath has the same
bound}\label{the-genuine-induced-marginal-heat-bath-has-the-same-bound}}

These error estimates also hold for the genuine unit-rate heat bath of
the actual separated marginal, one of the operators denoted \(B_X\) in
the historical retained-mean reduction. Write \(\nu_X\) for the actual
marginal, including its entire density
\(e^{-\beta S_{\mathrm{ext}}}\prod_{e\in E}Z_e\) relative to retained
Haar. For a retained link \(b\), let \(\overline K_b\) condition in this
very law on all other retained links, and set

\[
 B_X=\sum_{b\notin E}(I-\overline K_b).
 \tag{7.1a}
\]

Each \(e_c\) reads only its fixed observed collar. In a stationary
\(\nu_X\) one-link refresh it therefore changes only when that collar
reads the updated link; both endpoint fourth and second moments are the
actual ones already proved. The same bounded-incidence argument gives

\[
 \langle E_f,B_XE_f\rangle\le C\beta^{-1}M_f,
 \qquad
 \|1_{[\delta,\infty)}(B_X)E_f\|_2^2
 \le C(\beta\delta)^{-1}M_f.
 \tag{7.1b}
\]

This is not obtained by resetting the density to Wilson. The kernel
\(\overline K_b\) resamples using every induced normalizer, and need not
equal the original one-link kernel or its compressed update.

There is an exact current duality for this genuine marginal heat bath as
well. Eliminated current roots satisfy \(P_XT_e=0\). For retained \(b\),
\(t_b=P_XT_b\) obeys
\(\overline K_b t_b=\mathbb E(T_b\mid X\setminus\{b\})=0\), by the
original one-link score identity and the tower property, and
\(\|t_b\|_2\le C\). Thus

\[
 |\langle v,P_X\tau_g\rangle|
 \le C\sqrt{M_g\langle v,B_Xv\rangle}.
 \tag{7.1c}
\]

For spectral projections of this specific \(B_X\), (7.1b)--(7.1c) also
give the nonlinear-error/current cross bound
\(C\beta^{-1/2}\sqrt{M_fM_g}\). The ordinary error/quadratic cross term
is still unbounded. No identification with a sharply selected
\(\mathcal B_*\) heat bath or a mixed-history native operator is
asserted.

There is also an actual Dirichlet-form comparison. Put

\[
 M_f^{\rm ret}=\sum_cf_c(m_c-Em_c),\qquad
 Q_f^{\rm act}=\sum_cf_c(q_c-Eq_c).
\]

Their difference is E\_f exactly. Locality and the L2 bounds imply
\(\mathcal E(Q_f^{\rm act})\le CM_f\). Hence

\[
 \left|\mathcal E(M_f^{\rm ret})
             -\mathcal E(Q_f^{\rm act})\right|
 \le C\beta^{-1/2}M_f.
\tag{7.2}
\]

The analogous polarized bound is C beta\^{}-1/2 sqrt(M\_f M\_g). This
compares weighted Dirichlet covariances, not ordinary variances.

\hypertarget{an-additional-actual-negative-order-current-estimate}{%
\subsection{An additional actual negative-order current
estimate}\label{an-additional-actual-negative-order-current-estimate}}

For completeness, the exact score duality supplies a useful bound that
does not require the Taylor argument. For deterministic original weights
g, let

\[
 \tau_g=\tfrac13\sum_eg_eT_e,\qquad
 T_e=\Delta_eS_0-\beta|\nabla_eS_0|^2.
\]

Reference~\cite{CT} gives K\_eT\_e=0 and \textbar\textbar T\_e\textbar\textbar\_2
\(\leq\) C. Thus, for every G in L2,

\[
 |\langle G,\tau_g\rangle|
 \le C\sqrt{M_g}\,\mathcal E_{\rm orig}(G)^{1/2}.
\tag{8.1}
\]

For any deterministic original-law observation B, use its exact
compressed A\_B and set t\_g=P\_B tau\_g. Restricted to range(P\_B),
(8.1) implies by the spectral theorem or Riesz duality

\[
 \boxed{\quad
 \|A_B^{-1/2}t_g\|_2^2\le CM_g,
 \qquad
 \|1_{(0,\varepsilon]}(A_B)t_g\|_2^2\le C\varepsilon M_g.
 \quad}
\tag{8.2}
\]

The inverse is taken on positive spectrum. Equation (8.1) itself proves
that t\_g is orthogonal to the kernel; no current kernel mode is removed
by assumption.

For a spectral projection Q of A\_X, combining (8.1) with the small
error energy gives

\[
 |\langle QE_f,Q(P_X\tau_g)\rangle|
 \le C\beta^{-1/2}\sqrt{M_fM_g}.
\tag{8.3}
\]

These estimates control actual current interactions with the nonlinear
error. They do not control the ordinary error's interaction with the
shared ordinary retained mean.

\hypertarget{the-exact-failed-upgrade}{%
\subsection{The exact failed upgrade}\label{the-exact-failed-upgrade}}

The exact centered source decomposition is

\[
 P_XR_f-ER_f=Q_f^{\rm act}+E_f.
\tag{9.1}
\]

The left side is the actual retained ordinary source, and both
right-side variables are under the same actual Wilson marginal. The
quadratic Q\_f\^{}act uses the local chart variables of Section 3. It is
not the Gaussian quadratic benchmark evaluated under covariance (beta
S)\^{}+.

The available local remainder bound yields

\[
 |\operatorname{Cov}(e_c,e_d)|
 \le\|e_c-Ee_c\|_2\|e_d-Ee_d\|_2
 \le C\beta^{-1},
\]

uniformly in c,d, but gives no decay with their separation. Summing
therefore gives exactly the squared-l1 count in (1.5). Define

\[
 N_{\rm eff}(f)=\frac{(\sum_c|f_c|)^2}{M_f}
\]

for a nonzero mark. Then

\[
 \operatorname{Var}(E_f)\le C\frac{N_{\rm eff}(f)}\beta M_f.
\tag{9.2}
\]

For a uniform mark N\_eff is the number of roots. This proves an
extensive higher-order remainder only when that count is at most a fixed
multiple of beta. It supplies no bound uniform in the full allowed
volume regime.

The small Dirichlet bound does not improve the missing low spectrum: it
controls \(\int\lambda\,d\nu_{E_f}(\lambda)\), whereas the required
variance is \(\int d\nu_{E_f}(\lambda)\). Large weight at arbitrarily
small lambda is compatible with (1.3). A source-specific negative-order
estimate for E\_f, or a spatial connected bound on its root covariance,
would be new input.

Finally, even a small bound on \textbar\textbar E\_f\textbar\textbar{}
alone would not justify deleting the projected cross term. For any later
common observation/projection P,

\[
 \|P(Q_f^{\rm act}+E_f)\|_2^2-\|PQ_f^{\rm act}\|_2^2
 =2\operatorname{Re}\langle PQ_f^{\rm act},PE_f\rangle
     +\|PE_f\|_2^2.
\tag{9.3}
\]

The ordinary cross term is not the current cross term bounded in (8.3).
The true compact-jet comparison consequently leaves TWO explicit
obligations: a covariance comparison for Q\_f\^{}act under its actual
law, and an extensive spatial/low-spectral estimate for the nonlinear
error and its ordinary cross term.

\hypertarget{normalizations-and-limits-of-applicability}{%
\subsection{Normalizations and limits of
applicability}\label{normalizations-and-limits-of-applicability}}

All estimates are dimensionless. The original physical squared factor
c\_beta\^{}-2 a\^{}-4 stays outside. Constants do not absorb original
parent masses, selected branch probabilities, saved source grades, age
factors, or coupling ratios. The proof uses deterministic original
marks; random history coefficients have their own field/geometry
responses.

The observed local gauges, cutoffs and Taylor polynomial are analytic
devices in an exact decomposition. They do not redefine the physical
source or replace its readout. Every actual source and denominator
response is in m\_c, H\_c and e\_c.~Full compact integration is used
throughout this theorem; the sharp-domain cusp and later induced-density
issues remain outside its local C3 premise.

The positive results are the actual local beta\^{}-1/2 remainder, its
beta\^{}-1 extensive jump energy, the high-spectral error bound C/(beta
delta), and the negative-order current estimate. None closes the actual
retained-ordinary low-spectral estimate (11), mixed-domain admission, or
S94.

\newpage

\hypertarget{part-vi.-exact-source-pressure-identities-and-a-positivity-obstruction}{%
\section{Part VI. Exact source-pressure identities and a positivity
obstruction}\label{part-vi.-exact-source-pressure-identities-and-a-positivity-obstruction}}

This part examines whether the actual
retained-ordinary mean estimate can be turned into a source-pressure or dual-Dirichlet
bound. It does not prove that missing bound. It gives a precise
obstruction to reusing the Wilson positive-convolution argument for a
positive ordinary-source tilt, keeps the complete compact readout and
scalar factors, and derives a weaker source-averaged susceptibility
estimate under the original multipoint moment bound.

\hypertarget{the-literal-complete-su2-one-link-kernel}{%
\subsection{The literal complete SU(2) one-link
kernel}\label{the-literal-complete-su2-one-link-kernel}}

Use the original Wilson normalization, and let only one original link
differ from the identity. Write its quaternion as

\[
 u=(\cos\theta,\sin\theta\,n),\qquad |n|=1.
\]

The six incident plaquettes give exactly

\[
 S_0(u)-S_0(I)=12(1-\cos\theta).
 \tag{T1}
\]

For the COMPLETE uniform original directional-square source, all
affected roots must be counted. The central root contributes
\(12c_X^2\sin^2\theta\). The other eighteen incident roots each
contribute one nontrivial directional-force square, for a further
\(18c_X^2\sin^2\theta\). On a torus large enough that this fixed stencil
has no periodic self-identification,

\[
 D(u)=30c_X^2\sin^2\theta.
 \tag{T2}
\]

The full Casimir response obeys \(\ell=4C_F(N\#\mathcal P-S_0)\), with
\(C_F=3c_X^2/4\) in this convention. Therefore

\[
 -\frac13\{\ell(u)-\ell(I)\}
 =12c_X^2(1-\cos\theta).
 \tag{T3}
\]

Thus the COMPLETE ordinary source variation is

\[
 R(u)-R(I)
 =30c_X^2\beta\sin^2\theta
   +12c_X^2(1-\cos\theta).
 \tag{T4}
\]

These coefficients agree with a separate
complete-source calculation of the induced layer. The eighteen
neighboring roots are included; this is not the central-root
approximation. The entire field-independent value \(R(I)\) is retained.

For a positive real ordinary source parameter \(\eta>0\), the resulting
raw compact-star kernel is

\[
 k_\eta(u)
 =\exp\{-\beta[S_0(u)-S_0(I)]+\eta R(u)\}
 =C_\eta\exp\{A\cos\theta+B\sin^2\theta\},
 \tag{T5}
\]

where

\[
 A=12\beta-12\eta c_X^2,
 \qquad B=30\eta c_X^2\beta>0,
 \qquad C_\eta=e^{\eta R(I)-12\beta+12\eta c_X^2}>0.
 \tag{T6}
\]

Any additional original exterior constant multiplies \(C_\eta\) and
remains present. A positive scalar normalization cannot alter
positive-definiteness.

\hypertarget{positive-real-weights-do-not-give-a-positive-convolution-operator}{%
\subsection{Positive real weights do not give a positive convolution
operator}\label{positive-real-weights-do-not-give-a-positive-convolution-operator}}

\hypertarget{proposition}{%
\subsubsection{Proposition}\label{proposition}}

For every finite real \(A\) and every \(B>0\), the function

\[
 f(\theta)=\exp\{A\cos\theta+B\sin^2\theta\}
 \tag{T7}
\]

is not positive definite on \(U(1)\). Consequently the literal central
\(SU(2)\) kernel in (T5) is not the kernel of a positive convolution
operator.

\hypertarget{proof}{%
\subsubsection{Proof}\label{proof}}

Discard only the positive scalar factor \(e^{B/2}\), which is restored
without changing the conclusion. The Laurent extension of the remaining
function is

\[
 F(z)=\exp\left\{
 \frac A2(z+z^{-1})-\frac B4(z^2+z^{-2})
 \right\},\qquad z\in\mathbb C\setminus\{0\}.
 \tag{T8}
\]

This is holomorphic on the punctured plane. Its Laurent coefficients
\(c_m\) are precisely the Fourier coefficients of \(F(e^{i\theta})\),
and the Laurent series converges absolutely on every circle \(|z|=r\),
\(0<r<\infty\).

If \(f\) were positive definite on \(U(1)\), all Fourier coefficients
\(c_m\) would be nonnegative. Absolute convergence would then imply

\[
 |F(ir)|\le\sum_{m\in\mathbb Z}c_m r^m=F(r)
 \qquad(r>0).
 \tag{T9}
\]

Direct evaluation gives

\[
 \log|F(ir)|-\log F(r)
 =\frac B2(r^2+r^{-2})-\frac A2(r+r^{-1}).
 \tag{T10}
\]

The right side is positive for all sufficiently large \(r\), since
\(B>0\), contradicting (T9). Hence at least one Fourier coefficient is
negative.

A continuous positive-definite function on \(SU(2)\) restricts to a
positive-definite function on every subgroup. Restricting (T5) to the
diagonal \(U(1)\) subgroup produces (T7). Therefore the \(SU(2)\)
function is not positive definite either. For a smooth central function
on a compact group, positivity of the convolution operator is equivalent
to nonnegative representation Fourier coefficients, and hence to
positive-definiteness. This proves the asserted convolution failure.

The elementary Fourier implication used above can also be seen directly
on \(U(1)\): convolution acts on the Fourier mode \(e^{im\theta}\) by
its coefficient \(c_m\), so positivity requires every \(c_m\ge0\).

\hypertarget{exact-scope-of-this-obstruction}{%
\subsection{Exact scope of this
obstruction}\label{exact-scope-of-this-obstruction}}

The original Wilson character expansion has nonnegative coefficients and
supplies the positive transfer operators used in the arbitrary-set
moment theorem. Equation (T10) shows that this same positivity cannot be
invoked merely because the positive ordinary-source Boltzmann factor is
pointwise positive. The directional-square tilt has already destroyed
positive-definiteness in a literal compact one-link kernel with the
COMPLETE ordinary readout and full Casimir response.

This is an obstruction to that particular proposed extension of the
Wilson convolution/trace-Hölder argument. It does not establish that
every possible enlarged transfer construction fails. It does not show
that an ordinary-source susceptibility bound is false, and it says
nothing against the Yang--Mills continuum construction. It also does not
exclude a different one-sided source argument: the proved failure
concerns positive \(\eta\), for which \(B>0\).

\hypertarget{the-exact-global-pressure-identity-retains-the-ordinary-mean}{%
\subsection{The exact global pressure identity retains the ordinary
mean}\label{the-exact-global-pressure-identity-retains-the-ordinary-mean}}

Let \(F\) be the unchanged complete ordinary source,
\(P=\mathbb E(\cdot\mid\mathcal B)\) any actual completed positive
observation, and use its actual base marginal law \(\nu\). Define the
normalized conditional source logarithm

\[
 K(\eta;b)
 =\log\mathbb E(e^{\eta F}\mid\mathcal B=b),
 \qquad
 \Psi(\eta)=\log\mathbb E e^{\eta F}
 =\log\mathbb E_\nu e^{K(\eta;\mathcal B)}.
 \tag{T11}
\]

All fiber normalizers are retained in this formula: if \(Z_b(\eta)\) is
the raw conditional source integral, then
\(K(\eta;b)=\log Z_b(\eta)-\log Z_b(0)\), while \(Z_b(0)\) remains in
the actual marginal \(\nu\). One does not drop it from the outer
integration.

Differentiating at source zero gives the exact identity

\[
 \boxed{
 \Psi''(0)
 =\mathbb E_\nu K''(0;\mathcal B)
   +\operatorname{Var}_\nu K'(0;\mathcal B)
 =\mathbb E\operatorname{Var}(F\mid\mathcal B)
   +\operatorname{Var}(PF).}
 \tag{T12}
\]

The first term is controlled by the complete residual estimates. The
second is the actual retained-ordinary-mean problem. A local zero-free
disc for every conditional compact normalizer controls conditional
cumulants; it does not remove this second term or establish a zero-free
disc for the global source partition function.

For example, the conditional functions \(e^{n\eta}\) and \(e^{-n\eta}\)
are entire and zero free, while their equally weighted positive mixture
is \(\cosh(n\eta)\), whose closest complex zeros have distance
\(\pi/(2n)\) from the origin and whose logarithmic curvature at zero is
\(n^2\). This is a counterexample to an inference from conditional
zero-free factors to a global zero-free disc. It is not asserted to be a
Wilson counterexample.

For the original full physical source prescription, primitive
higher-grade contact, current, reference and boundary partners must be
added in their actual prescribed grades. Equation (T12) is the exact
ordinary exponential-source direction, not an identification of those
other grades with ordinary conditional cumulants. If the selected
domains themselves move with a source, their original surface and
selector derivatives also remain.

\hypertarget{a-limited-actual-law-source-averaged-susceptibility-bound}{%
\subsection{A limited actual-law source-averaged susceptibility
bound}\label{a-limited-actual-law-source-averaged-susceptibility-bound}}

The arbitrary-set original moments do yield a quantitative pressure
estimate, but it does not select the original point \(\eta=0\).

Let \(F_f=\sum_cf_cr_c^0\), \(b_f=\max_c|f_c|>0\), and
\(L_f=\sum_c|f_c|\). The complete local source satisfies the pointwise
estimate

\[
 |r_c^0|\le C\left(1+\beta\sum_{p\in\mathcal L(c)}s_p\right),
 \tag{T13}
\]

where \(\mathcal L(c)\) is its original fixed determining record and all
incidences are bounded. For arbitrary weights \(0\le w_p\le1\),
convexity of the log moment generating function and layer-cake
interpolation of the weights give, from the arbitrary-set estimate,

\[
 \log\mathbb E\exp\left(\frac{t\beta}{6}\sum_pw_ps_p\right)
 \le\frac{\log K_t}{6}\sum_pw_p.
 \tag{T14}
\]

For completeness, \(w_p=\int_0^1 1_{\{w_p\ge u\}}du\). Jensen's
inequality for the convex log moment functional bounds its value at
\(w\) by the integral of its values on these indicator vectors; apply
the proved arbitrary-set estimate to each vector.

Applying (T14) to the bounded-overlap plaquette weights generated by
(T13), choose a fixed \(\delta>0\) such that

\[
 \log\mathbb E\exp\left(\frac\delta{b_f}|F_f|\right)
 \le C\frac{L_f}{b_f}.
 \tag{T15}
\]

Constants are independent of volume and beta in the stated original
moment regime. In particular \(|\mathbb EF_f|\le CL_f\).

Set

\[
 \psi_f(\eta)=\log\mathbb E
            e^{\eta(F_f-\mathbb EF_f)}.
\]

For real \(|\eta|\le\delta/b_f\), Jensen gives \(\psi_f(\eta)\ge0\),
while (T15) gives \(\psi_f(\eta)\le C L_f/b_f\). Convexity therefore
bounds the derivatives at the half-radius points by

\[
 -CL_f\le\psi_f'(-\delta/(2b_f))
 \le\psi_f'(\delta/(2b_f))\le CL_f.
 \tag{T16}
\]

For the EXACT source-tilted original law

\[
 d\mu_\eta=\frac{e^{\eta F_f}}{\mathbb E e^{\eta F_f}}\,d\mu,
\]

one has \(\psi_f''(\eta)=\operatorname{Var}_{\mu_\eta}(F_f)\). Thus

\[
 \boxed{
 \int_{-\delta/(2b_f)}^{\delta/(2b_f)}
       \operatorname{Var}_{\mu_\eta}(F_f)\,d\eta
 \le CL_f.}
 \tag{T17}
\]

The average over this interval is at most \(Cb_fL_f\). For uniform marks
of magnitude one on \(m\) roots, this is \(Cm=CM_f\). For general marks
it is the displayed \(b_fL_f\) count; it is not silently replaced by
\(M_f\).

This bound concerns a source average over the exact tilted family, with
all original scalar normalizers. The target uses \(\mu_0=\mu\).
Convexity permits a narrow peak of \(\psi_f''\) centered at zero, so
(T17) does not prove the target susceptibility at zero. No tilted law is
substituted for the original law in the construction.

\hypertarget{where-the-source-specific-dual-dirichlet-attempt-stops}{%
\subsection{Where the source-specific dual-Dirichlet argument
stops}\label{where-the-source-specific-dual-dirichlet-attempt-stops}}

The complete current enjoys an exact dual estimate because each \(T_e\)
is conditionally centered under the actual original one-link kernel:

\[
 \langle v,P\tau_f\rangle
 =\frac13\sum_e\widetilde f_e
       \langle(I-K_e)v,T_e\rangle.
 \tag{T18}
\]

For a directional action \(S_{e,\mu}\), the actual compact
integration-by-parts identity instead reads

\[
 \mathbb E\left[
  \{\Delta_eS_{e,\mu}-\beta F_e\cdot F_{e,\mu}\}v\right]
 =-\mathbb E\langle F_{e,\mu},\nabla_ev\rangle,
 \qquad F_e=\sum_\nu F_{e,\nu}.
 \tag{T19}
\]

Its force product contains all cross directions. Replacing it by
\(\beta|F_{e,\mu}|^2\) is the unsupported step. A pointwise
factorization \(F_e\cdot V_e=\sum_\mu|F_{e,\mu}|^2\) can fail even at
the conditional minimum of a good star. In the actual \(SU(2)\) staple
convention, take

\[
 h_1=2(\cos\alpha\,e_0+\sin\alpha\,e_1),\quad
 h_2=2(\cos\alpha\,e_0-\sin\alpha\,e_1),\quad h_3=2e_0.
\]

Each pair consists of two actual unit staples, realizable by the
distinct exterior links. Their total is \((4\cos\alpha+2)e_0\), so the
conditional minimum is \(u=e_0\) and the total force there is zero. The
first two individual directional forces are nonzero and opposite for
\(\alpha\ne0\). Every incident plaquette is good for sufficiently small
\(|\alpha|\). Thus no finite vector \(V_e\) can satisfy the proposed
factorization at that point, and the expression
\(V_e=(\sum_\mu|F_{e,\mu}|^2)F_e/|F_e|^2\) is singular. This observation
restricts that local construction; it does not preclude a different
global source-specific representation using additional fields or
derivatives.

Likewise, the proved retained-source Dirichlet estimate bounds

\[
 \int\lambda\,d\sigma_{r_f}(\lambda)\le CM_f,
 \qquad r_f=P_BR_f-\mathbb ER_f.
 \tag{T20}
\]

It does not bound \(\int_{(0,1]}d\sigma_{r_f}\). The compact flat-jet
estimate supplies high spatial derivatives at a particular flat bath,
whereas \(\sigma_{r_f}\) is the spectral measure of the full actual
configuration-space dynamics. No proven intertwining identifies those
two spectral decompositions over the nonlinear marginal law.
Substituting the flat spatial estimate into (T20) as if they were the
same spectral cutoff is therefore another invalid step.

The remaining problem is an actual-law estimate on the retained ordinary
spectral mass or its annealed covariance. The current score identity,
local source zero-free discs, generalized multipoint moments, and the
source-averaged bound above do not yet provide it.

\begin{thebibliography}{9}
\bibitem{CT} \emph{Wilson current concentration and conditional Gaussian fluctuations}, \texttt{lattice\_theory/block\_and\_transfer/conditional\_estimates/wilson\_current\_concentration\_and\_conditional\_gaussian\_fluctuations.tex} in this collection. Part~A, subsection ``A direct proof of the extensive bridge'' (the separated-link observation $X$, the error $\epsilon_h$ and $P_X$).
\bibitem{OS} \emph{Ordinary spectral estimates and mixed current transfer}, \texttt{lattice\_theory/block\_and\_transfer/conditional\_estimates/ordinary\_spectral\_estimates\_and\_mixed\_current\_transfer.tex} in this collection.
\end{thebibliography}

\newpage
\end{document}
